
\documentclass[11pt]{article}

\usepackage{amsfonts}
\usepackage{amsmath}
\usepackage{amscd}
\usepackage{amssymb}
\usepackage{ifthen}
\usepackage{indentfirst}
\usepackage{natbib}
\usepackage{rotating}
\usepackage[colorlinks=true,allcolors=blue]{hyperref}
\usepackage{url}
\usepackage{doi}
\usepackage{tikz-cd}

\DeclareMathOperator{\idfun}{id}
\DeclareMathOperator{\isContr}{isContr}
\DeclareMathOperator{\isProp}{isProp}
\DeclareMathOperator{\isSet}{isSet}
\DeclareMathOperator{\dom}{dom}
\DeclareMathOperator{\cod}{cod}
\DeclareMathOperator{\qinv}{qinv}
\DeclareMathOperator{\binv}{biinv}
\DeclareMathOperator{\isEquiv}{isequiv}
\DeclareMathOperator{\refl}{refl}
\DeclareMathOperator{\inl}{inl}
\DeclareMathOperator{\inr}{inr}
\DeclareMathOperator{\successor}{succ}
\DeclareMathOperator{\inductor}{ind}
\DeclareMathOperator{\recursor}{rec}

\newcommand{\sloop}{\text{loop}}
\newcommand{\sbase}{\text{base}}
\newcommand{\ssurf}{\text{surf}}

\newcommand{\set}[1]{\{\,#1\,\}}
\newcommand{\proptrunc}[1]{\lVert #1 \rVert}
\newcommand{\termtrunc}[1]{\lvert #1 \rvert}
\newcommand{\bigproptrunc}[1]{\left\lVert #1 \right\rVert}

\newcommand{\ints}{\mathbb{Z}}
\newcommand{\nats}{\mathbb{N}}
\newcommand{\real}{\mathbb{R}}
\newcommand{\universe}{\mathcal{U}}
\newcommand{\rats}{\mathbb{Q}}

\newcommand{\fatdot}{\,\cdot\,}

\let\emptyset=\varnothing

\newcommand{\REVISED}{\begin{center} \LARGE REVISED DOWN TO HERE \end{center}}
\newcommand{\MOVED}[1][equation]{\begin{center} [#1 moved] \end{center}}

\begin{document}

\title{On Generalizations of Equality, Logic, and Set Theory
    in Mathematics of the Past Century}
\author{Charles J. Geyer}
\maketitle

\tableofcontents

\section{License and Source Code}

This work is licensed under a Creative Commons
Attribution-ShareAlike 4.0 International License
(\url{http://creativecommons.org/licenses/by-sa/4.0/}).

The \LaTeX\ source for this document is
at \url{https://github.com/cjgeyer/Survey}.

\pagebreak[3]
\section{Introduction}

\begin{quotation}
In brief, a philosophy of Mathematics is not convincing unless it is
founded on an examination of Mathematics itself.
Wittgenstein (and other philosophers) have failed in this regard.
\\
\hspace*{\fill} --- \citet{maclane}
\end{quotation}

This document started with two questions raised by Roy Cook
in the foundations interest
group (FIG) of the Minnesota Center for Philosophy of Science, one March 19,
2026 and one much earlier (back in fall 2025).
\begin{itemize}
\item What is really a substructure?
\item What is finite, and can you define the natural
    numbers when you don't even know what finite is?
\end{itemize}

The point of this document is to explain how contemporary mathematics
makes these
questions uninteresting.  They do not even arise.  Those brainwashed in
logic and traditional
set theory cannot understand what this is about without learning
newer ideas from abstract algebra, category theory, constructive analysis,
and type theory, especially homotopy type theory.

So this document evolved into a history of mathematics since 1900
with special attention to these subjects.

Special attention is paid to several themes.
\begin{itemize}
\item We do not want equality in the logic.
\item We do not want classical logic of any kind.  Intuitionistic logic
    is used in many areas of contemporary mathematics.  If we want
    classicality principles like the law of excluded middle or the axiom
    of choice, add them to theorem statements as explicit assumptions
    rather than have them in the logic by default.
\item We actually want many logics: categorical logics, type theories,
    lambda-calculi, and even classical logic.  We are not interested
    in worshiping one true logic.
\item Type theory replaces logic and set theory.  It does not need them.
\item Functions and, more generally, morphisms of categories are far
    more important than any other concept in mathematics.
    They are fundamental, not part-whole relations.
\end{itemize}

\section{Up to Isomorphism}
\label{sec:up-to-isomorphism}

Long ago in FIG we read \citet{mazur}, which explains that ``equal''
in contemporary mathematics has been demoted from logical equality
to ``isomorphic''.\footnote{\citeauthor{mazur} also presents an extreme
structuralist view that the Yoneda lemma of category theory means that
isomorphism is determined by the morphisms of objects to other objects.}

How many trivial groups (having exactly one element) are there?
The up-to-isomorphism view says just one.
If you insist that there are many, then all are isomorphic.
Contemporary mathematics says that arguing about whether the only element of a
trivial group is ``really'' the empty set or a banana or Julius Caesar
is just silly.

And the same for everything else in mathematics, equality up to isomorphism
is good enough.
I know about Frege's Caesar problem.  But I don't care, and neither does
any contemporary mathematician.
Of course, even structuralist philosophers agree with this position, more
or less (of course there are many different structuralist views).

As we shall see (Section~\ref{sec:homotopy-equivalence} below),
homotopy type theory says that even isomorphism is too strong and
needs to be replaced with a weaker concept: homotopy equivalence.

\section{The Elementary Theory of the Category of Sets}

Set theory was reconceptualized in the 1960's by
Lawvere.  His elementary theory of the category of sets (ETCS)
\citep{lawvere-etcs}
has a way of making the issue of substructure disappear.

Long ago in FIG we read \citet{leinster}, which is an elementary
explication of ETCS.  ETCS is motivated by category theory and topos theory,
but is actually just a theory in first-order logic like ZFC
(Zermelo–Fraenkel set theory with the axiom of choice).  ZFC is a one-sorted
theory with variables ranging over sets and one non-logical binary relation
symbol $\in$ read as membership.  ETCS is a two-sorted
theory with some variables ranging over sets and some variables ranging over
functions and one non-logical binary relation
symbol $\circ$ read as composition of functions.\footnote{``Ranging over''
is sloppy terminology.  We just have variables in formulas.  If there is only
one sort of variable, we don't have to say what those variables refer to.
We could do set theory without ever mentioning ``sets,'' although this
goes against the way human languages work.
When there are two sorts of variables,
we do have to keep them separate, but do not have to name them either.}
We can (this is a matter of taste) also include non-logical constants
$\idfun_A$ for every set $A$, which are identity functions.
If we do not include these constants,
instead we include axioms that say identity functions exist and
are unique.\footnote{It is also possible to state ETCS as a one-sorted
theory with variables ranging over functions and one ternary relation
which holds for functions $f$, $g$, and $h$ if and only if $h = g \circ f$.
Then there are some axioms that
say that some functions behave like identity functions (they obey the identity
laws of a category, equations \eqref{eq:identity} and \eqref{eq:identity-too}
below) and then, of course, each identity function can be associated with
a set (object).  But in this functions-only view, there are no sets.  The
only role of sets in ETCS or objects in categories (Section~\ref{sec:category}
below) is to indicate which functions in ETCS or morphisms in categories
are composable.  \citet[p.~9]{maclane-working}
gives details.\label{foot:arrows-only}}

In short, in ZFC sets and set membership are the fundamental concepts,
whereas in ETCS sets and
functions and function composition are the fundamental concepts, and the
set membership and subset relations must be defined in terms of them.

To define membership \citep[Section~1]{leinster}
we first define \emph{terminal} sets.
These are sets $T$ such there is at exactly one function $A \to T$ for any
set $A$.  An axiom of ETCS asserts that terminal sets exist.
Pick one and denote it $1$ (it turns out that, as in category theory,
all terminal sets are isomorphic,\footnote{Proof: Suppose $T$ and $T'$ are
terminal.  Then the functions $T \to T' \to T$ and $T' \to T \to T'$ are
unique, hence must be the identity functions $\idfun_T$ and $\idfun_{T'}$.
But this says the unique functions $T \to T'$ and $T' \to T$ are
inverse functions of each other, hence isomorphisms.
\label{foot:terminal}
See Section~\ref{sec:category} below for the definition of isomorphism.}
so it does not matter which one we pick).

Then ETCS defines an element $x$ of a set $A$ to be a function $x : 1 \to A$.
So worrying about whether (or exactly how) $x$ is in $A$ is a non-starter.
Functions (including $x$) aren't in sets, rather they go between sets
(from one to another).

You may be wondering what function evaluation now means.  What does $f(x)$
denote if $x$ is not an element of $A$ in the sense of ZFC?  It denotes
the function composition $f \circ x$, that is, the composition of functions
$$
\begin{CD}
   1 @>x>> A @>f>> B
\end{CD}
$$
This all makes sense because by definition the composite arrow $f \circ x$
maps $1 \to B$ hence is an element of $B$.

To define the subset relation \citep[Discussion of Axiom~8]{leinster}
we first define \emph{injective} functions.
In ETCS we say a function is \emph{injective} if
it has the left cancellation property: if $x$ and $y$ are elements of
the domain of $f$, then $f \circ x = f \circ y$ implies $x = y$.
Of course, this is just the usual definition in different notation.

In ETCS a \emph{subset} is just an injective function into a set.\footnote{
Of course, an element is also an injective function, trivially so, because
$1$ has only one point (proof: there is one function $1 \to 1$ because $1$
is terminal).  Thus ETCS does not distinguish between an element $x$ and
the singleton set $\{x\}$ which is a subset.}
Hence it makes no philosophical sense to ask in what way this function
has some sort of part-whole relationship to its codomain.

\citet{leinster} following \citet{lawvere-etcs} argues that mathematicians
outside of pure set theory \citep[the subject of][]{jech} only use ETCS,
which is a weaker theory than ZCF
\citep[Section~``How strong are the axioms?'']{leinster}, lacking any analogue
of the ZFC axiom schema of replacement.
It can even be argued that mathematicians outside of pure set theory mostly
use even less, only naive set theory \citep{halmos}.

ETCS does feel very different from ZFC.  \citet{leinster} does not mention
arbitrary union and intersection, and notes that even arbitrary disjoint
unions are not guaranteed by the ETCS axioms.  Moreover, \citet{leinster}
only mentions the power set once, saying that in ETCS one uses instead
the function space $2^X$, the set of all functions $X \to 2$,
instead of the power set of $X$.  Nevertheless, the function concept is
powerful enough to do all of mathematics.

\section{Abstract Algebra}
\label{sec:abstract-algebra}

The up-to-isomorphism view goes back to the rise of abstract algebra in
the 1920's in the work of Noether and coworkers
\citep[Chapter~6]{kleiner} and in the textbooks of
\citet{van-der-waerden} and \citet{birkhoff-maclane}.

Consider abstract group theory, where the only objects of interest are groups
and the only functions of interest are group homomorphisms.
The first isomorphism theorem
\citep[of group theory,\footnote{Almost every algebraic theory
has a similar theorem, even set theory
\citep[Theorem~I.2.7]{aluffi}.}][Theorem~II.8.1]{aluffi}
says, if $f : G \to H$ is a group homomorphism,
then it may be decomposed as in the following commutative diagram
(more on commutative diagrams in Section~\ref{sec:diagram} below)
\begin{equation} \label{diag:isomorphism}
\begin{tikzcd}
   G \rar{f} \dar[swap]{q} & H \\
   G / \ker(f) \rar{\tilde{f}} & f(G) \uar[swap]{i}
\end{tikzcd}
\end{equation}
where $G / \ker(f)$ is the group $G$ with equality redefined to mean
$x =_{G / \ker(f)} y$ if and only if $f(x) =_H f(y)$,\footnote{The
notation $=_A$ means equality as defined for the object $A$, much
more on this in what follows, particularly Sections~\ref{sec:constructive}
and~\ref{sec:identity} below.}
where $q$ is the quotient map $x \mapsto x$
(with equality defined differently on the domain and codomain)
where $\tilde{f}$ is the canonical isomorphism
(the same as $f$ except with equality defined differently on the domain
of $f$ and the domain of $\tilde{f}$ and with the codomains being as
indicated in the diagram),
and where $i$  is the canonical injection $x \mapsto x$.

The theorem has zero interest in whether $f(G)$ and $H$ have some
sort of part-whole relationship (or exactly what kind if so).  We do have
$f(G)$ a subset of $H$ in the sense of ETCS because the canonical injection
is an injection.  But much more importantly, the fact that this arrow
indicates a group homomorphism says that $f(G)$ is a subgroup of $H$
rather than a mere subset.
Similarly $G / \ker(f)$ and $f(G)$ are isomorphic, hence
each is a subset of the other in the sense of ETCS.  And that is all group
theorists really need (or should want).\footnote{
The notation $\ker(f)$ denotes a subgroup of $G$ and
the notation $f(G)$ a subgroup of $H$ \citep[Section~II.6]{aluffi},
but the notation $G / \ker(f)$ does
not denote a subgroup of anything in diagram~\ref{diag:isomorphism}.
Rather, $G / \ker(f)$ is a factor group of $G$, that is,
$\ker(f) \times (G / \ker(f))$ is isomorphic to $G$
\citep[Section~II.3.4]{aluffi}.}

A group $G$ is \emph{simple} if for any homomorphism $f : G \to H$
the group $G / \ker(f)$ as described by the first isomorphism theorem
is isomorphic to either $G$ itself or the trivial group.
Simple groups are like prime factors of groups
(the Jordan-H\"{o}lder theorem, \citealp[Theorem~IV.3.1]{aluffi}).

A great accomplishment of twentieth century mathematics was the
classification of the finite simple groups, finished in 1980 or perhaps 2004
\citep{aschbacher}.
It is perhaps the most amazing example of mathematics delivering amazing
complexity that was not built in at the start.
To even describe the classification (all of the finite simple groups)
takes a large book
\citep[more than 300 pages if front matter is included]{wilson}.
The whole proof to this day (as far as I can see) has not been digested
and written up in a straightforward way, although work on this is ongoing.

So the notion of structure-preserving function (called homomorphism in group
theory) subsumes and generalizes the notion of subset or even
structure-preserving subset (here subgroup).
And this way of thinking is 100 years old and has gotten stronger and stronger
over the years.

Another moral of the story is that isomorphism is not particularly interesting.
If $f$ in the first isomorphism theorem is itself an isomorphism, then all
four groups in diagram~\ref{diag:isomorphism} are isomorphic, and we haven't
learned anything.  It is only when $f$ is a homomorphism that is not an
isomorphism and $G / \ker(f)$ is not isomorphic to $G$ or the trivial group
that we learn something (get a non-trivial factorization).
So we are interested in all the homomorphisms, not just the isomorphisms.

\section{Constructive Analysis}
\label{sec:constructive}

It may seem strange that we defined a quotient group in the preceding section
by redefining equality (rather than using equivalence classes), but this
is commonplace in constructive mathematics.  It goes back at least to
\citet{bishop} and is also seen in Martin-L\"{o}f type theory (MLTT)
\citep{martin-lof-mltt}, on which
homotopy type theory (HoTT)
\citep[hereinafter referred to as the HoTT book]{hott-book} is based.
\defcitealias{hott-book}{the HoTT book}

\Citet{bishop} defines a set to be rules for producing elements of the set
and rules for saying which elements of the set are equal.  So equality is
a defined relation, which must be an equivalence relation (reflexive,
symmetric, and transitive) and may be defined differently for each set.

In HoTT, if we want to be pedantically correct,
we write $=_A$ rather than just $=$ to mean equality for the type $A$.

So we did not need to think of equivalence classes at all (much less to
think of them as sets) to define a quotient group.  The explanation gets
much more complicated if based on equality-in-the-logic and set theory.

In constructive mathematics, which is mathematics that does
not use the law of excluded middle (LEM) anywhere, putting equality
in the logic no longer works well.  LEM says that for any proposition $P$,
either $P$ or not $P$, that is, $\forall P, P \vee \neg P$.
A constructively provable equivalent to LEM is double negation
removal $\forall P, \neg \neg P \to P$.
When applied to negated propositions, LEM and double negation removal
are constructively valid (provable without LEM), that is,
\begin{align*}
   & \forall P, (\neg P) \vee (\neg \neg P)
   \\
   & \forall P, (\neg \neg \neg P) \to (\neg P)
\end{align*}
are both rules of constructive mathematics.  Thus if we want to talk about
$x$ and $y$, the propositions
\begin{align*}
   & x = y
   \\
   & \neg (x = y)
   \\
   & \neg \neg (x = y)
\end{align*}
are all different.
The result of this mess is that to get real math done, you had better
make sure equals means exactly what you want it to mean and ignore the
mess.  And this is what Bishop did.  (Much more on this theme in all
that follows.)

\section{Category Theory}
\label{sec:category}

We have already seen how in ETCS the set membership and subset relations are
defined in terms of functions.  In category theory, functions are generalized
to morphisms, which are like any abstraction: morphisms have some properties of
functions but not all properties.

\subsection{Definition and Axioms}

A category (like ETCS) has two sorts of variables, those ranging over
\emph{objects} and those ranging over \emph{morphisms}.\footnote{One can also
make category theory a one-sorted theory with variables ranging over morphisms,
as discussed in footnote~\ref{foot:arrows-only} above.}
And it has a non-logical binary function
symbol $\circ$ for composition of morphisms.  And (according to taste, like
in ETCS) it has
a non-logical constant symbol $\idfun_A$ for every object $A$,
which are identity morphisms.

Every morphism $f : A \to B$ goes from one object to another.\footnote{So
if we are being very fussy, we can add to the formalism functions $\dom$
and $\cod$ each taking morphisms to objects such that $f : A \to B$ means
the same as $A = \dom(f)$ and $B = \cod(f)$.}
Like for functions, we can call $A$ the domain and $B$ the codomain of $f$.
But if we like to emphasize that morphisms are different from functions,
we can call $A$ the source and $B$ the target.
Morphisms are sometimes called arrows.

The axioms for category theory are very sparse.  For a diagram
\begin{equation} \label{diag:abc}
\begin{CD}
   A @>f>> B @>g>> C
\end{CD}
\end{equation}
there is a morphism $g \circ f : A \to C$.\footnote{That is,
$g \circ f$ exists if and only if $\dom(g) = \cod(f)$.}  And for a diagram
$$
\begin{CD}
   A @>f>> B @>g>> C @>h>> D
\end{CD}
$$
we have the associative law $h \circ (g \circ f) = (h \circ g) \circ f$.
And we also have the identity laws that say for any $f : A \to B$ we have
\begin{subequations}
\begin{align}
    f \circ \idfun_A & = f
    \label{eq:identity}
    \\
    \idfun_B \circ f & = f
    \label{eq:identity-too}
\end{align}
\end{subequations}
Identities are unique.\footnote{If $\idfun_A$ and $\idfun_A'$ were
both identities for object $A$, then the identity laws imply
$\idfun_A = \idfun_A \circ \idfun_A' = \idfun_A'$.\label{foot:identity}}

\subsection{Diagram Chasing}
\label{sec:diagram}

Category theory makes much use of \emph{commutative diagrams}.
These are diagrams of directed graphs in which the nodes of the graph
are objects of the category and the edges of the graph are morphisms
of the category.  We say the diagram is commutative if parallel paths
are equal.

For example, to say that the following diagram is commutative
\begin{equation}
\label{diag:composition}
\begin{tikzcd}
   A \rar[swap]{f} \arrow[rr, bend left, "g \circ f"] & B \rar[swap]{g} & C
\end{tikzcd}
\end{equation}
is to express the existence of composite morphisms.
For another example, to say that the following diagram is commutative
\begin{equation}
\label{diag:associative}
\begin{tikzcd}
    A \arrow{d}[swap]{f} \arrow{dr}{g \circ f} \\
    B \arrow{r}{g} \arrow{dr}[swap]{h \circ g} & C \arrow{d}[swap]{h} \\
    & D
\end{tikzcd}
\end{equation}
is to express the associative law: that all three paths from $A$ to $D$
are equal is to say that
$h \circ (g \circ f) = h \circ g \circ f = (h \circ g) \circ f$.

For another example, to say that the following diagram is commutative
\begin{equation}
\label{diag:identity}
\begin{tikzcd}
    A \arrow{r}{f} \arrow{d}[swap]{\idfun_A} \arrow{dr}{f} &
    B \arrow{d}{\idfun_B} \\
    A \arrow{r}{f} & B
\end{tikzcd}
\end{equation}
is to express the identity laws.

It is hard to describe how useful this concept is
without going through some complicated applications.
Commutative diagrams, like \eqref{diag:isomorphism} above,
organize whole proofs.  The diagram itself is (almost) the proof.
You just have to check that all the details it asserts are correct.

\subsection{Categories in Which Morphisms are Functions}

The axioms of ETCS say that ETCS is a category with sets as the objects
and functions as the morphisms.  But ETCS needs more axioms than just
the category axioms to make its
objects and morphisms act like sets and functions.

We note, if it isn't already obvious, that the objects of a category need
not be sets and do not have to have elements and that the morphisms of a
category need not be functions, and there need be no sense in which they
can be evaluated (examples below).

But for many categories, the objects are sets with extra structure and the
morphisms are structure-preserving functions.\footnote{To say that morphisms
are actually functions means that composition of morphisms is actually
composition of functions, that is, $h = g \circ f$ means $h(x) = g(f(x))$
for all $x$ in the domain of $f$, which is also the domain of $h$.
But, in general, composition of morphisms only has to satisfy
the laws; it does not need to be defined in terms of something else.}
Some of these are shown in Table~\ref{tab:categories}
\ifthenelse{\equal{\arabic{page}}{\pageref{tab:categories}}}
{(this page)}{(p.~\pageref{tab:categories})}.
\begin{table}
\caption{Categories of Structured Sets and Structure-Preserving Functions}
\begin{center}
\begin{tabular}{lp{2.75in}}
objects & morphisms \\
\hline
sets & functions \\
groups & group homomorphisms \\
rings & ring homomorphisms \\
fields & field homomorphisms \\
vector spaces & linear functions \\
topological spaces & continuous functions \\
metric spaces & uniformly continuous functions \\
topological groups & continuous group homomorphisms \\
smooth manifolds & differentiable functions \\
Lie groups & differentiable group homomorphisms \\
topological vector spaces & continuous linear functions \\
partially ordered sets & monotone functions \\
measure spaces & measurable functions \\
statistical models & functions measurable in the random variable
    and continuous in the parameter
\end{tabular}
\end{center}
\label{tab:categories}
\end{table}

\subsection{Isomorphisms}
\label{sec:iso}

In category theory, a morphism $f : A \to B$ is \emph{iso},
that is, $f$ is an \emph{isomorphism}, if there exists $g : B \to A$ such that
\begin{subequations}
\begin{align}
   g \circ f & = \idfun_A
   \label{eq:isodef-a}
   \\
   f \circ g & = \idfun_B
   \label{eq:isodef-b}
\end{align}
\end{subequations}
In all of the categories in Table~\ref{tab:categories} this category theoretic
notion of isomorphism coincides with the pre-existing notion of isomorphism
that existed before category theory was invented.

When $f$ is an isomorphism, we write a $g$ satisfying \eqref{eq:isodef-a}
and \eqref{eq:isodef-b} as $f^{-1}$ and say that $f^{-1}$ is the \emph{inverse}
of $f$.  This is OK because inverses are unique if they exist.\footnote{Suppose
\eqref{eq:isodef-a} and \eqref{eq:isodef-b} and also hold with $g$ replaced
by $g'$.
Then $g = g \circ \idfun_B = g \circ (f \circ g') = (g \circ f) \circ g'
= \idfun_A \circ g' = g'$.
\label{foot:inverse}}

Isomorphism answers an important philosophical question.  What is a
topological property?  Or a group theoretic property?  Or a
whatever-area-of-mathematics-we-are-discussing property?
The answer is a property that does not distinguish between isomorphic
objects of the category of those objects and their morphisms.
Admittedly this is rather tautological.  If you are interested in a property
that distinguishes isomorphic topological spaces, then you are not doing
topology but rather some other area of mathematics.
This explains the up-to-isomorphism view (Sections~\ref{sec:up-to-isomorphism}
and~\ref{sec:abstract-algebra} above and this section).

Interestingly, this does not work for ZFC-like set theory.
Sets are isomorphic if they have the same cardinality,
but set theory is interested in lots of properties that distinguish
sets of the same cardinality.

\subsection{Preorders}
\label{sec:preorder}

A \emph{preorder} is a category in which for any two objects $A$ and $B$
there is at most one morphism $A \to B$.  Since there is one arrow $A \to A$,
the identity morphism, this is the only arrow $A \to A$.
Composition of morphisms says if there are arrows $A \to B \to C$, then there
is an arrow $A \to C$.  Since the only issue is whether there is
an arrow between any two objects,
this looks like a relation (a binary predicate).
And the identity and composition laws say this is a reflexive and transitive
relation.  Such a relation is called a \emph{preorder}.
So to make this look more like conventional
mathematics we write $A \le B$ rather than $A \to B$.

Objects $A$ and $B$ in this category are isomorphic
if $A \le B$ and $B \le A$.\footnote{Proof: going back to arrow notation,
there can be only one arrow
$A \to B \to A$ or $B \to A \to B$, so these must be identity arrows;
hence $A \to B$ and $B \to A$ are inverse arrows.
This is essentially the same proof as in footnote~\ref{foot:terminal}.}
Orders do not allow isomorphic but unequal objects.
So a preorder becomes a partial order if
we add the assumption that isomorphic objects are equal
($A \le B$ and $B \le A$ implies $A = B$).

And a partial order becomes a total order if we add the assumption that
for all objects $A$ and $B$, either $A \le B$ or $B \le A$.

The moral of the story is that morphisms in a category do not have to
be anything like functions, and categories do not have to be anything
like those in Table~\ref{tab:categories}

\subsection{Groupoids}
\label{sec:groupoid}

A \emph{groupoid} is a category in which every morphism is iso.
A \emph{group} is a groupoid with only one object.

If a group has only one object, call it $*$, how is that a group?
The morphisms of the category satisfy the group axioms.
Composition of morphisms is the group operation (called addition or
multiplication according to taste).  This operation is associative
by the associativity law for categories,
there is one identity morphism $\idfun_{\textstyle *}$
that behaves correctly because of the identity laws for categories.
Moreover, every morphism has an inverse by the assumption that all
morphisms are iso.

A group is thus a category in which the objects are uninteresting;
the arrows are everything.  But the arrows need be nothing like functions.
This view of groups may seem strange, but was actually well known in
group theory before category theory was invented.
A \emph{permutation group} is exactly this construction when the object
is a finite set and the morphisms are all iso (invertible functions
from this finite set to itself).\footnote{An invertible function from
a set of chairs to itself tells the occupants of the chairs where to
go to permute themselves.  If sitting in chair $x$, then move to chair $f(x)$.}
An \emph{automorphism group} is exactly this construction when the object
is any mathematical object whatsoever and the morphisms are isomorphisms
in any sense that agrees with the definition in category theory:
in any category having an object $A$, the set of all isomorphisms $A \to A$
forms a group.

The idea of group presented in this section is different from the one
in Table~\ref{tab:categories}.  Here the category is just one group.
There it is all groups and the morphisms are group homomorphisms.
Here we seem to have lost the group homomorphisms, but we can get them
back (next section).

\subsection{Functors}

A \emph{functor} between categories is a function that maps
objects to objects and morphisms to morphisms and respects the laws,
that is, if $F : \mathcal{C} \to \mathcal{D}$ is a functor
and we have diagram \eqref{diag:abc}
in the source category $\mathcal{C}$ then we have the diagram
\begin{equation}
\label{diag:functor}
\begin{CD}
   F(A) @>{F(f)}>> F(B) @>{F(g)}>> F(C)
\end{CD}
\end{equation}
in the target category $\mathcal{D}$. That is,
$$
   F(g \circ f) = F(g) \circ F(f)
$$
Then we can prove the functor also maps
identities to identities\footnote{In \eqref{diag:functor} replacing
one or the other of the functions by an identity function gives
for $f : A \to B$ in the source category
$F(f) \circ F(\idfun_A) = F(f)$ and
$F(\idfun_B) \circ F(f) = F(f)$.  This proves $F(\idfun_A)$ is an
identity morphism for object $F(A)$ and similarly with $A$ replaced by $B$.
This uses uniqueness of identities
(footnote~\ref{foot:identity}).\label{foot:functor-identity}}
\begin{equation} \label{eq:functorality-identity}
   F(\idfun_A) = \idfun_{F(A)}
\end{equation}
and inverses (when they exist) to inverses\footnote{In
diagram~\eqref{diag:functor} suppose $f$ and $g$ are inverse functions.
then $A = C$ and $F(A) = F(C)$ so $F(f)$ and $F(g)$ are also inverse functions
(because the composite arrow is an identity).
This uses functorality of identities \eqref{eq:functorality-identity}
and uniqueness of inverses
(footnote~\ref{foot:inverse}).\label{foot:functor-inverse}}
$$
   F(f^{-1}) = F(f)^{-1}
$$

A functor between categories that are groups
(in the sense of Section~\ref{sec:groupoid})
is a group homomorphism in the sense of group theory.\footnote{Before
category theory, a group homomorphism was defined to be
a function $f : G \to H$ between groups that respected the group
multiplication $f(x y) = f(x) f(y)$ (if we denote the group operation by
juxtaposition).  But that is just what \eqref{diag:functor} says.
Then we have to prove that a group homomorphism respects identities
and inverses, but that comes from the functor definition
(footnotes~\ref{foot:functor-identity} and~\ref{foot:functor-inverse}).}

A functor between categories that are preorders (in the sense of
Section~\ref{sec:preorder} is a monotone function in the sense of the
theory of orders and preorders.\footnote{This is obvious.
Diagram \eqref{diag:functor}, translated into the symbology that replaces
$\to$ with $\le$ says $x \le y$ implies $f(x) \le f(y)$.}

\subsection{Homotopy}
\label{sec:homotopy}

Category theory was created \citep{eilenberg-maclane}
to do algebraic topology and homotopy theory.
What are they?

Algebraic topology is the study of topology using algebraic methods,
mostly using groups associated with topological spaces.
A homotopy is a continuous deformation of a continuous function.
Originally, a homotopy between functions $X \to Y$ was defined to be
a (jointly) continuous
function $I \times X \to Y$ where $I$ denotes the unit interval
of the real number system $\set{x \in \real : 0 \le x \le 1}$.
Here $I \times X$ denotes the product topological space, which has the product
topology.
The idea is that we have two functions $f_0 = f(0, \fatdot)$
and $f_1 = f(1, \fatdot)$
and more generally $f_t = f(t, \fatdot)$, and $f_0$ and $f_1$ are
continuous functions $X \to Y$ (morphisms of the category of topological
spaces and continuous functions) and $f_t$ changes continuously from $f_0$
to $f_1$ as $t$ goes from 0 to 1.

But this particular construction is irrelevant.  The actual concept is
continuous deformation.  When there is a homotopy like the one above we
say that $f_0$ and $f_1$ are \emph{homotopic} or \emph{homotopy equivalent}.

This leads to the definition of \emph{homotopy category} in which the
objects are topological spaces and the morphisms are equivalence classes
of homotopy equivalent functions.
Alternatively, as in constructive mathematics, we can redefine what equals
means rather than use equivalence classes.  Morphisms are equal if they
are homotopic.

Isomorphisms in this category are the same as in any category.  A morphism
$f : A \to B$ is iso if there is a morphism $g : B \to A$ such that
\eqref{eq:isodef-a} and \eqref{eq:isodef-b} hold but now equals means
homotopic to, that is, $g \circ f$ is homotopy equivalent to the identity
function on $A$ and $f \circ g$ is homotopy equivalent to the identity
function on $B$.

When objects of this category (topological spaces) are isomorphic,
we say they are \emph{homotopy equivalent}.

One moral of the story is here is another category in which the morphisms
are not functions.\footnote{Or they are functions with a defined notion
of equality (homotopy).}

Now for the algebraic topology.  A \emph{path} in a topological space $A$
is a continuous function $I \to A$.  The relation of there being a path
from a point $x$ to a point $y$ is an equivalence relation (reflexive,
symmetric, and transitive).\footnote{Proof: the constant path takes $x$ to $x$,
hence reflexive.  Every path $p$ has an inverse $p^{-1}$ defined by
$p^{-1}(t) = p(1 - t)$, $0 \le t \le 1$, hence symmetric.  Every two paths $p$
and $q$ such that the endpoint of $p$ is the startpoint of $q$ have a
path concatenation $r = p \cdot q$
defined by $r(t) = p(2 t)$, $0 \le t \le 1/2$
and $r(t) = q(2 t - 1)$, $1 / 2 \le t \le 1$, hence transitive.}

For any $x \in A$, we say the set of all paths from $x$ to $x$ is the
\emph{fundamental group} of $A$ for \emph{basepoint} $x$.  How is this
a group?  Path concatenation is the binary operation.  Every path from
$x$ to $x$ is then iso by the symmetry property.  This is the
notion of group-as-category as in Section~\ref{sec:groupoid}.

Two topological spaces that are topologically isomorphic (the term
for this used in topology before category theory was invented was
\emph{homeomorphic}) must have the same fundamental group but not
vice versa.  Thus the property of having the same fundamental group
is a weaker property than isomorphism.
Why would anyone want a weaker property?  Because it says something
different and perhaps useful and perhaps easier to calculate.

The trouble with equivalent concepts is the law of conservation of
mathematical difficulty.  Duality often does not help.
But a weaker concept perhaps can.
This is related to the notion that morphisms that are not iso can be
more important than isomorphisms (Section~\ref{sec:abstract-algebra} above).

But this document is not about algebraic topology, so we leave the subject
here.  (It will come back in Section~\ref{sec:higher-induction} below.)

We do want to say something else about homotopy.
Since homotopies are continuous functions, we can make homotopies
of homotopies, jointly continuous functions $I \times I \times X \to Y$.
And homotopies of homotopies of homotopies,
jointly continuous functions $I \times I \times I \times X \to Y$.
And so on and on with arbitrarily many ``homotopies of.''
This theme will return at the end of Section~\ref{sec:identity}.

\subsection{Category Theory as a Foundation of Mathematics}

\citet{lawvere-thesis,lawvere-cat-of-cat}
also proposed category theory as a foundation of mathematics.
The foundational theory is the category of all categories
in which the objects are categories and the morphisms are functors.

Of course, if the category of all categories had itself as one of the objects
we could reproduce Russell's paradox.
So like in NBG (von Neumann--Bernays--G\"{o}del) set theory, we do not allow
that.  We say that the objects of the category of all categories are
\emph{small categories} and that does not include the category of all
categories (which thus is analogous to the universe of NBG set theory,
which is a proper class).

Of course the category of all categories contains the category of sets and
functions.  Whether we take these to be ZFC sets and functions or NBG sets
and functions or ETCS sets and functions or HoTT sets and functions
\citepalias[Chapter~10]{hott-book} is a matter of taste.
But it also contains the rest of mathematics.

\subsection{Categorical Logic}
\label{sec:categorical-logic}

I preface this section with an admission that I am incompetent to write
on the subject, having never used it.  So I should not have the temerity
to write about it.  But I cannot just skip it either.  It is part of the
history.  But this section may have mistakes.

The simplest part of this story is that categorical logic generalizes
model theory (the branch of logic so called).  Theories are categories
(that's another part of the story) and models are functors from theories
into any other category.  There is no reason to have models in the category
of sets and functions (like classical model theory does).
In classical logic, models provide meaning (semantics).
Thus this feature of categorical logic is called \emph{functorial semantics}.

This is very powerful.  If one looks for models of the theory of groups in
the category of topological spaces and continuous functions, one discovers
topological groups.  If one looks in the category of smooth manifolds
and differentiable functions, one discovers Lie groups.

The next part of the story is that logical theories themselves also
become categories.  The objects are well-formed formulas and the morphisms
are logical implications.\footnote{This is where I am making a mistake.
We define objects (formulas) to be equal if they are the same except for
variable renaming.}  Thus logical theories are categories that are preorders
in the sense of Section~\ref{sec:preorder} above.\footnote{But they have
more properties than being merely preorders.}
Actually, this was already required for functorial semantics to make sense:
a functor goes from one category, here a logical theory, to another,
here a logical model.

The last part of the story is that certain categories, called toposes or
topoi (for those who like Greek plurals for Greek-derived terms
(singular: topos)),
have what is called an internal logic.  They have \emph{subobject classifiers},
also called \emph{truth value objects} that set up a correspondence between
subsets-as-injective-morphisms (like in ETCS) and maps from the set
(not subset) to the subobject classifier.  In set theory (including ETCS)
we can think of this as the correspondence between subsets $A \to S$ and
maps $S \to \Omega$, where for classical logic the subobject classifier
$\Omega$ is isomorphic to
the two-valued set $\{\text{true}, \text{false}\}$.  But other toposes have
other internal logics.  Most are constructive, so the subobject classifier
is a Heyting algebra.

A lot more could be said on this subject \citep{goldblatt}, but we
leave this subject here.

\subsection{Applied Category Theory}

It can be argued that category theory has more applications in the rest
of mathematics than set theory.  No doubt set theory, at least naive
set theory \citep{halmos}, is fundamental to several areas of
mathematics,\footnote{ In point set topology,
the fundamental object is a topological space, which
consists of a set $X$ and a family $\mathcal{O}$ of subsets of $X$,
which is called a \emph{topology} and whose elements
are called the \emph{open sets}, satisfying some axioms: the whole space
and the empty set are open sets, and arbitrary unions and finite intersections
of open sets are open.  Similarly, in probability theory, a fundamental object
is a measurable space, which
consists of a set $X$ and a family $\mathcal{A}$ of subsets of $X$,
which is called a \emph{sigma-algebra} and whose elements
are called the \emph{measurable sets}, satisfying Kolmogorov's axioms,
and then a \emph{probability measure} is a function $\mathcal{A} \to \real$
satisfying another of Kolmogorov's axioms.\label{foot:sigma-algebra}}
although some of these areas can be replaced by more category-theoretic
approaches \citep{johnstone}.
By the 1950's set theory had percolated down to middle-school mathematics,
where it was called ``new math.''

In the second half of the twentieth century, category theory percolated
through advanced mathematics, at least what may be called ``naive category
theory'' by analogy with ``naive set theory'' even though there is no book
with the former title.  As examples we can cite \citet{aluffi} which is
full of category theory, commutative diagrams, and diagram-chasing proofs.
And one might reply, sure, but that's algebra, what category theory was
designed for.  But look at \citet{lang} and \citet{beattie-butzmann}, which
do real and functional analysis, which you might think should be all
for every $\varepsilon > 0$ there is a $\delta > 0$ such that $\ldots,$
but which have some of that but also a lot of algebra so these books are
also full of commutative diagrams and diagram-chasing proofs.
Moreover, they are full of definitions, concepts, and theorems
that are categorically
inspired.  \citet[Preface]{beattie-butzmann} even explain that much of the
research literature they are digesting uses deep results from category theory
which they have left out of the book so as not to bother their audience.

But category theory can also be used to redo logic
(Section~\ref{sec:categorical-logic} above)
and computer programming.  One can think of a pure functional programming
language (such as Haskell or OCaml or the languages written to implement
proof assistants, such as ML, Rocq, and Agda) as a computer language having
only one feature: functions.  All you can do is define functions and apply
them.  These functions are the morphisms of the category, and composition
is following one with another, which is the same as plugging in the results
of one into the other (as its argument).  So category theory influences the
design of these computer languages.\footnote{It is a joke but also a
technically correct mathematical statement that a \emph{monad}
on a category is a monoid object in the monoidal category of endofunctors.
A joke because it mystifies more than explains.  But monads are a big deal
in functional programming.}

As far as I know category theory has only filtered down to advanced
undergraduate mathematics (courses mostly for math majors).\footnote{Nowadays,
courses below this level are largely free of set theory too.}  It is not
clear to me that the K-12 education establishment even wants to repeat
the new math debacle with category theory.

\section{Homotopy Type Theory}

For those brainwashed in logic and ZFC-like set theory, HoTT is really strange.
It has no need for logic and set theory.  It eats them.  It has logic
and set theory inside it.  They become parts of type theory.
So type theory is not a theory in first-order logic or any other logic.

Also MLTT has no axioms (formulas in some logic asserted to be true).
Instead it is formulated as rules for constructing types.
HoTT does have one axiom: it adds to MLTT the univalence axiom
(Section~\ref{sec:univalence} below).

\subsection{Types}

Types, the fundamental concept of type theory, are not defined
extensionally.
Instead types are intensional.  Type theory, according to
\citet[near the beginning of the first lecture]{bauer-hott},
is a theory of constructions.  Types are characterized by
rules for producing terms of the type.\footnote{This is like Gentzen's
natural deduction, just a lot of rules.  The rules for type theory
are also often presented like in natural deduction: a horizontal line
with premises above and the conclusion the rule provides below,
as in Appendix~A.2 of \citetalias{hott-book} and throughout \citet{rijke}.
Appendix~A.1 of \citetalias{hott-book} formalizes MLTT in Church's
simply typed lambda-calculus.}
This sounds like Bishop's definition of set, but is more general.

In MLTT and HoTT the fundamental judgment is $a : A$
read as $a$ is a term of type $A$.\footnote{\citet{ladyman-presnell} say
it is more philosophically correct to say ``token'' rather than ``term''
here.}$^,$\footnote{There are two other \emph{judgments} in the
metatheory rather than in HoTT.  These are $A$ is a type, and types or
terms are \emph{definitionally (or judgmentally) equal}
\citepalias[Section~1.1]{hott-book}.}
There are three different interpretations of $a : A$ used throughout HoTT
in addition to the ``neutral'' one just stated.
We can say $a$ is an element of the set $A$, or we can say $a$ is a point
in the space $A$ (which says the rules of $A$ give it additional structure
over and above being a mere set),
or we can say $A$ is a proposition (interpreted
as true if it is inhabited and false otherwise) and $a$ is a proof of $A$.
These readings are summarized in Table~\ref{tab:term-type}
\ifthenelse{\equal{\arabic{page}}{\pageref{tab:term-type}}}
{(this page)}{(p.~\pageref{tab:term-type})}.
\begin{table}
\begin{center}
\begin{tabular}{ll}
term & type \\
\hline
proof & proposition \\
element & set \\
point & space
\end{tabular}
\end{center}
\caption{Readings of $a : A$.  Alternatives are token instead of term and
witness or evidence or certificate instead of proof.}
\label{tab:term-type}
\end{table}
 
The reading about proofs of propositions shows us that type theory
also eats proof theory.  To prove a
proposition $A$ is the same as constructing a term of type $A$.  So truth
becomes completely unproblematic: now defined to be types
being inhabited.

Some think it misleading when $p : P$ to call $p$ a proof of $P$.
Those people say $p$ is
a \emph{witness} to or \emph{evidence} of or \emph{certificate} for
the truth of $P$.  But in MLTT
and HoTT the truth of $P$ does not mean anything other than that $p : P$ for
some $p$, that is, does not mean anything other than that $P$ is inhabited.
Hence an inhabitant of $P$ proves $P$.

It is important to understand that these readings are mere analogies.
They are not part of HoTT.  There are only terms and types, no matter
how you talk about them.

\subsection{Identity Types}
\label{sec:identity}

Also like in \citet{bishop}, as mentioned above, equality is a defined
relation for each type.
If $A$ is a type, and $x, y : A$, then we write $x =_A y$ to say that
$x$ and $y$ are equal terms \emph{according to the way equality is defined
for type $A$.}  But MLTT and HoTT go beyond that in considering this equality\
to be itself a type, denoted $(x =_A y)$ or $\text{Id}_A(x, y)$.
And, like any type, this type can have terms.

So consider $p, q : (x =_A y)$.
What are $p$ and $q$?  Different (perhaps) ways $x$ and $y$ can be equal.
This caused a lot of confusion about MLTT until HoTT arose.
When we are thinking homotopically, we interpret $p$ and $q$ as paths
from $x$ to $y$ (as in algebraic topology and homotopy theory).  Then
it is no surprise at all that there can be more than one path from $x$
to $y$.  Similarly, if we are thinking that $x$ and $y$ are proofs of
$A$ considered as a proposition, it is no surprise that there can be
different proofs and maybe different ways proofs can be logically equivalent.
It is only when we say that there can be different ways that elements
of sets can be equal that this sounds bizarre to those brainwashed in
logic and ZFC-like set theory.

We should stress that, unlike in algebraic topology and homotopy theory,
these paths $p, q : (x =_A y)$ are not defined in terms of some other concepts.
Specifically, they are not defined as continuous functions $I \to A$
where $I$ is the
unit interval of the real line.  These are synthetic rather than analytic.
They are just always there for any types.  And they have the properties
given them by the rules of the type and no other properties.\footnote{Identity
types also have rules (Section~\ref{sec:path-induction} below).  So there
are rules not only for type $A$ but also for type $(x =_A y)$.}

Sometimes, like in the rest of math, we will omit the subscript on the
equals sign, writing $x = y$ instead of $x =_A y$, when the type
of $x$ and $y$ is obvious from the context.  But this does not indicate
a different concept.  In MLTT and HoTT $x = y$ only makes sense
when $x$ and $y$
have the same type.  If that type is $A$, then it means $x =_A y$.
In MLTT and HoTT it makes no sense to ask whether terms of different types are
equal.\footnote{So like in ETCS and category theory, set membership
and the subset relation have to be defined otherwise than in ZFC-like
set theories.  No sets-as-types can have a subset-superset relationship
(or any kind of part-whole relationship).}

Also, like in the rest of math, equality is an equivalence relation:
symmetric, reflexive, and transitive.  Reflexivity means $x = x$ for
any term $x$ of any type.  And what the preceding sentence means in HoTT
is that $(x = x)$ is an inhabited type (truth is inhabitation).  So for
any $x : A$ there exists
$$
   \refl_x : (x =_A x)
$$
which we can think of as the constant path from $x$ to $x$,\footnote{If
it were a function $f : I \to A$, where $I$ is the closed unit interval
of the real line, then it would be the function defined by $f(t) = x$,
$0 \le t \le 1$.} but, as always in HoTT, this is synthetic rather than
analytic, just there rather than defined in terms of something else.

Transitivity requires that $x, y, z : A$ and $x = y$ and $y = z$
imply $x = z$.  Propositions-as-types (Section~\ref{sec:curry-howard} below)
requires we write this as the function type (logical implication is a special
case of function)
$$
   ((x = y) \times (y = z)) \to (x = z)
$$
then the notion of currying says we can write a function of two variables
as a function of one variable whose value is a function of one variable,
that is, write $f(x)(y)$ instead of $f(x, y)$ so the type
becomes\footnote{To be really pedantic (a set theorist is a person who
thinks all functions have only one argument)
a function $f : A \times B \to C$ has arguments
that are ordered pairs so should be written $f\bigl((x, y)\bigr)$, which is
really ugly, so no one does that.  The default is that a type
$A \to B \to C$ is read as right associative $A \to (B \to C)$.
So $f : A \to (B \to C)$ is a function that maps $a \mapsto g$ where
$a : A$ and $g : B \to C$.  So $f(a) = g$ and $f(a)(b) = g(b)$ is
a point in $C$.  So the notation $f(x, y)$ is really
ambiguous about whether the type
is $A \times B \to C$ or $A \to B \to C$.
The term currying for this comes from the name of the logician
Haskell Curry.\label{foot:currying}}
$$
   (x = y) \to (y = z) \to (x = z)
$$
\citepalias[Lemma~2.1.2 and the recursor for product types,
equation (1.5.2)]{hott-book}.  This function of two variables, which are
paths $p : x = y$ and $q : y = z$ is called \emph{path concatentation}
and written $p \cdot q : (x = z)$.
Again, we can think of this as the path that is obtained by following $p$
and then following $q$, but it is really synthetic rather than analytic:
transitivity requires it is just there.

Symmetry requires that $x, y : A$ and $x = y$ implies
$$
   (x = y) \to (y = x)
$$
and this function, which can also be thought of as logical implication,
is written $p \mapsto p^{-1}$.
Again, we can think of $p^{-1}$ as $p$ travelled backwards,
but it is really synthetic rather than analytic.

It is a theorem that $\refl_x^{-1} \equiv \refl_x$
\citepalias[Lemma~2.1.1]{hott-book} ($\equiv$ means equal by definition
rather than just $=_A$ for some type $A$), which gives us another reason
for thinking of it as a constant path.  Also it acts like an identity
for path concatentation $p : x = y$ implies
$p = p \cdot \refl_y = \refl_x \cdot p$ and also all of the other identity
properties of a category
\citepalias[Lemma~2.1.4]{hott-book}.
Thus a type $A$ is a groupoid (Section~\ref{sec:groupoid} above)
when considered as a category with objects being terms $x : A$ and
paths $p : (x =_A y)$ being considered as morphisms $x \to y$
\citepalias[Section~2.1]{hott-book}.

All of the above may seem rather simple until we realize that any
statement about equality of paths involves the existence of paths between
paths (homotopies in homotopy theory, Section~\ref{sec:homotopy} above,
although, as usual, synthetic rather than analytic), that is, if
$p, q : (x =_A y)$, then we can wonder whether $p = q$ but that is
a question about whether the type $(p =_{(x =_A y)} q)$ is inhabited,
and that is a type of paths between paths (homotopies).
And then if that type has inhabitants $f$ and $g$, we can wonder whether
$f = g$ but then that is a question about whether the type
$f =_{\left(p =_{(x =_A y)} q\right)} g$ is inhabited, so a question about paths
between homotopies (homotopies of homotopies
or paths between paths between paths) and so forth \emph{ad infinitum}
(just like in homotopy theory, Section~\ref{sec:homotopy} above).

The story of identity types continues in Sections~\ref{sec:path-induction}
and~\ref{sec:higher-induction} below.

\subsection{Propositions and Sets}

A type $A$ is a proposition \citepalias[Section~3.3]{hott-book} if it
has at most one path component, that is,
\begin{equation} \label{eq:proposition}
   \isProp(A) :\equiv \prod_{x, y : A} (x =_A y)
\end{equation}
where $:\equiv$ indicates a definition and the $\Pi$-type can be read
as either a Cartesian product or as logical for-all ($\forall$).
Here it is best read as the latter.  A type $A$ is a proposition
(\citetalias{hott-book} often says \emph{mere} proposition
to emphasize meaning this concept) if any two points are equal.
So in a sense it has at most one point, although the homotopical interpretation
is at most one path component (every two points are path connected).
This makes mere propositions quite boring.
The only interesting thing you can say
about them is whether they are inhabited (true) or not (false).

A type $A$ is a set \citepalias[Section~3.1]{hott-book} if every path component
is contractible, that is,
\begin{equation} \label{eq:set}
   \isSet(A) :\equiv \prod_{x, y : A} \prod_{p, q : (x =_A y)}
   (p =_{(x =_A y)} q)
\end{equation}
This is somewhat confusing, especially the notation, $p =_{(x =_A y)} q$,
which is the identity type $=_B$ where $B$ is itself an identity type
$(x =_A y)$.

In short, a type is a set when it is homotopically boring.  Points
are either in the same path component or they aren't and the path
components are uninteresting.

This may not seem much like ZFC, but
these HoTT sets behave much like ETCS sets \citepalias[Chapter~10]{hott-book}.
But ETCS has the axiom of choice
and HoTT doesn't (except if you add it as an extra assumption to prove
something, Section~\ref{sec:truncation} below),
so the HoTT sets make a constructive set theory.

An element of a set is, as we said above, just $a : A$, a term of a type.
But a subset is weird: in one sense just like ZFC and in
another sense wildly different.  It is a $\Sigma$-type
\citepalias[Section~3.5]{hott-book}
\begin{equation} \label{eq:subset}
   \sum_{x : A} P(x)
\end{equation}
where each $P(x)$ is a mere proposition.  In another notation from the rest
of mathematics, this means the same thing as $\set{x \in A: P(x)}$.
So HoTT subsets are $\set{x \in A: P(x)}$.  We do not reify them as
anything ``more fundamental.''

This is like ETCS in that subsets are not ``in'' the sets they are subsets
of in any sense.  But otherwise it is very different.  In ETCS (recall)
subsets are injective functions.  In HoTT they are dependent pair types.

A term of the type \eqref{eq:subset} is a pair $(x, p)$
where $x : A$ and $p : P(x)$.  Of course, where there is no proof (that is,
when $P(x)$ is not inhabited), there is no pair ether.
So \eqref{eq:subset} consists of pairs $(x, p)$ such that $x$ is a point
of the subset and $p$ is a proof that this $x$ is in the subset.

You may think it inelegant that we are carrying the proofs along with
the points.  But think again.  The alternative is to search all of mathematics
for the proof when you need it to continue another proof.  By carrying it
along the computer proof assistant does not have to search for $p$ when it
is needed.  This is called \emph{proof-relevant} logic
(more on this in Section~\ref{sec:truncation} below).

\subsection{The Curry-Howard Correspondence}
\label{sec:curry-howard}

\subsubsection{How Logic fits in to Type Theory}

MLTT and HoTT can seem horribly baroque compared to ZFC-like set theories.
So complicated!  They have lots of constructions and lots of rules
governing their constructions.  ZFC just has one non-logical symbol for the
membership relation and a few axioms.

But ZFC needs first order logic, and MLTT and HoTT do not.  So that is not
a fair comparison.  A good way to see that type theory is actually simpler
than logic plus set theory is to see that it does not need anything
not needed for them.  It just generalizes them.

So we copy Table~1 from \citetalias{hott-book} adding a bit from
Table~1.1 in \citet{rijke} to get our Table~\ref{tab:foo}
\ifthenelse{\equal{\arabic{page}}{\pageref{tab:foo}}}
{(this page)}{(p.~\pageref{tab:foo})},
which shows how logic fits into type theory,
the so-called Curry-Howard correspondence.
\begin{sidewaystable}
\begin{center}
\begin{tabular}{cllll}
Symbol & Types & Logic & Sets & Homotopy \\
\hline
$A$ & type & proposition & set & space \\
$a : A$ & term & proof & element & point \\
$B(x)$ & dependent type & predicate & family of sets & fibration \\
$b(x)$ & dependent term & conditional proof & family of elements & section \\
$0$ & empty type & $\bot$ & $\varnothing$ & empty space \\
$1$ & unit type & $\top$ & singleton set & one-point space \\
$A + B$ & sum type & $A \vee B$ & disjoint union & coproduct space \\
$A \times B$ & product type & $A \wedge B$ & set of pairs & product space \\
$A \to B$ & function type & $A \implies B$ & set of functions & function space
\\
$\sum_{x : A} B(x)$ & dependent pair type & $\exists_{x : A} B(x)$ &
    disjoint sum & total space \\
$\prod_{x : A} B(x)$ & dependent function type & $\forall_{x : A} B(x)$ &
    Cartesian product & space of sections \\
$\text{Id}_A$ or $=_A$ & identity type & equality (=) &
    $\set{(x, x) : x \in A}$ & path space \\
$A \to 0$ & & negation ($\neg A$) & &
\end{tabular}
\end{center}
\caption{Type Theory, Logic, Set Theory, and Homotopy Theory}
\label{tab:foo}
\end{sidewaystable}

\subsubsection{Negation, the Empty Type and Empty Function}

The last line of Table~\ref{tab:foo} does not really fit with the rest,
but we need it to fill out the Curry-Howard correspondence.
If $A$ is a type, its negation
is defined to be the type $A \to 0$, where $0$ denotes the empty type.
This function space is inhabited
if and only if $A$ is not inhabited.  In case $A \equiv 0$ the function
type $0 \to 0$ has exactly one term: the empty function having a vacuous
rule.\footnote{If $f : 0 \to 0$, then the rule
defines $f(x)$ for all $x : 0$ but there are no such $x$.
But there is such a function $f$
(all mathematicians agree).\label{foot:emptyfun}
Also, perhaps even more bizarrely, we can also call this the identity
function $0 \to 0$, as it does map $x \mapsto x$, but vacuously, there
being no such $x$.  This makes every type $A$ have an identity function
$A \to A$.
How do we formalize this?  In Section~1.7 \citetalias{hott-book}
defines the \emph{recursor} for the empty type to be exactly the rule
that for any type $A$ there is a function $0 \to A$, that is,
the type $0 \to A$ is inhabited.  Propositions-as-types interprets this
as the logical principle \emph{ex falso quodlibet}.}
Conversely, in case $A$ is inhabited, there can be no function $A \to 0$
because there is an $x : A$ but there is no point in the empty type for $f(x)$
to be.\footnote{More pedantically, the rule for construction of a function
$f : A \to B$ says you give a formula $f(x) :\equiv \Phi$, $x : A$, where
$\Phi$ is an expression in type theory that may have $x$ as a variable.
And, in order for this construction to be valid, you or a proof assistant
must prove $f(x) : B$ for all $x : A$.}

\subsubsection{Cartesian Products}

Another thing we should perhaps clarify in the table is that in type theory
a term of type $A \times B$ is an ordered pair $(a, b)$ where ordered pair
is a primitive concept of type theory.  In set theory, of course,
$$
   A \times B = \set{ (x, y) : x \in A \wedge y \in B }
$$
where $(a, b)$ is still an ordered pair but not a primitive concept
but rather something defined in terms of
more primitive concepts.\footnote{Such as the Kuratowski definition
$(a, b) = \{\{a\}, \{a, b\}\}$ where $\{a\}$ is the singleton set whose
only element is $a$ and $\{a, b\}$ is the unordered pair whose existence
is guaranteed by the pairing axiom of ZFC.  Of course, the existence
of singleton sets is also guaranteed by the pairing axiom plus the axiom of
extensionality because $\{a\} = \{a, a\}$.}

Yet another thing we should perhaps clarify in the table is that in type theory
a dependent function $f : \prod_{x : A} B(x)$ is just a function such
that $f(x) : B(x)$ so the codomain of the function depends on
the argument $x$.

In set theory, of course, arbitrary Cartesian
products are defined in exactly the same way,
if $\mathcal{A} = \set{ A_i : i \in I }$ is an
indexed family of sets then $\prod \mathcal{A} = \prod_{i \in I} A_i$
is the set of functions $f : I \to \bigcup \mathcal{A}$
satisfying $f(i) \in A_i$ for all $i$.
But in set theory, function is not primitive and must be defined in terms
of more primitive concepts.

One way to state the axiom of choice is
that if every element of $\mathcal{A}$ is nonempty, then $\prod \mathcal{A}$
is nonempty.  This includes the case where $\mathcal{A}$ is the empty family,
in which case the only inhabitant of $\prod \mathcal{A}$ is 
the empty function of footnote~\ref{foot:emptyfun}.

In HoTT \citepalias[Theorem~2.15.7]{hott-book} we do not need the axiom of
choice in order for a dependent product type to be inhabited.\footnote{We
still need to assume $B(x)$ is inhabited for
all $x : A$.  But then $\prod_{x : A} B(x)$ is inhabited without the axiom
of choice.  More pedantically, the rule for construction of a function
$f : \prod_{x : A} B(x)$ says you give a formula $f(x) :\equiv \Phi$,
$x : A$, where
$\Phi$ is an expression in type theory that may have $x$ as a variable.
And, in order for this construction to be valid, you or a proof assistant
must prove $f(x) : B(x)$ for all $x : A$.  So you do not need the axiom
of choice because the choices have already been made in the construction.}
This is sometimes called the type theory axiom of choice.
However, this does not have the all of the logical consequences of the
classical axiom of choice.  (More on this after some preliminaries.)

\subsubsection{Some Redundancy}

When the dependent type $B(a)$ does not
actually depend on $a$ we have
\begin{align*}
   \sum_{a : A} B & \equiv A \times B
   \\
   \prod_{a : A} B & \equiv A \to B
\end{align*}
So there is some redundancy in the basic HoTT operations.

\subsection{Propositional Truncation}
\label{sec:truncation}

If $A$ is a type, then there is another type $\proptrunc{A}$,
called the \emph{propositional truncation} of $A$, defined
as follows.  There is a function $\termtrunc{\fatdot} : A \to \proptrunc{A}$
(so for each $a : A$ there is $\termtrunc{a} : \proptrunc{A}$),
and for all $a, b : \proptrunc{A}$, we have $a = b$
(so equality in $\proptrunc{A}$ is defined so that all terms are equal).
This makes $\proptrunc{A}$ a mere proposition, regardless of what $A$ is.

In Table~\ref{tab:foo} we were a bit cavalier about two of the lines.
Various theorems in \citetalias{hott-book} prove that the operations take
mere propositions to mere propositions except for the sum types, which don't.
For example, the unit type $1$ is a mere proposition
but the sum type $1 + 1$ has two elements that
are not equal.  So it isn't a proposition.
So if we want to get what people usually
think of as logic, we have to squash the sums.  Logical or becomes
$\proptrunc{A + B}$ and the dependent sum type (logical exists) becomes
$\bigproptrunc{\sum_{x : A} B(x)}$
\citepalias[Definition~3.7.1]{hott-book}.

But \citetalias[Section~3.7]{hott-book}, also argues against the routine
use of propositional truncation.  If $P$, $Q$, and $R$ are mere propositions
and we want an implication $\proptrunc{P + Q} \to R$ it suffices to prove
$P + Q \to R$ because the recursion principle for propositional truncation
gives us the implication $\proptrunc{P + Q} \to R$.

What is the difference?  To prove $\proptrunc{P + Q}$, we prove $P$ or prove
$Q$ and forget which one.  To prove $P + Q$, we prove $P$ or prove
$Q$ and remember which one.\footnote{And how does the sum type remember?
A term of type $A + B$ has the form $\inl(a)$ for some $a : A$
or $\inr(b)$ for some $b : B$ \citepalias[Section~1.7]{hott-book}.
In the first case, it is part of $\inl(a)$ that $a : A$, and similarly
for the second case.}
Mathematically, there is no point to the forgetting.

And similarly for the dependent sum type, which we read as logical exists.
This is similar to the discussion of the subset relation \eqref{eq:subset}
above.
A term of $\sum_{x : A} P(x)$ is an ordered pair $(x, p)$ where
$x : A$ and $p$ is a proof of the proposition $P(x)$.  A term of its
propositional truncation forgets both $x$ and $p$.\footnote{More precisely,
for any $y : A$ and $q : P(y)$ we have $\termtrunc{(x, p)}
=_{\proptrunc{\sum_{x : A} P(x)}} \termtrunc{(y, q)}$.}
Mathematically, there is no point to the forgetting
(propositional truncation).\footnote{If you assume
$\proptrunc{\sum_{x : A} B(x)}$, then you know there exist $x : A$
and $y : B(x)$ but you do not know which ones.  So you can use $x$ and $y$
in your proof so long as what you are trying to prove is a mere proposition
that does not depend on the particular $x$ and $y$.
And how does the type system (the proof assistant) enforce this?  An
inhabitant of $\proptrunc{\sum_{x : A} B(x)}$ has the form
$\termtrunc{(x, y)}$ for $x : A$ and $y : B(x)$, where the vertical bars
indicate the propositional truncation, that is,
$\termtrunc{(x, y)} = \termtrunc{(x', y')}$ for any other such pair
$(x', y')$.  Hence the inverse of function extensionality,
equation \eqref{eq:happly} below, implies that any function (implication)
$f : \proptrunc{\sum_{x : A} B(x)} \to Q$ has to have
$f\bigl(\termtrunc{(x, y)}\bigr) = f\bigl(\termtrunc{(x', y')}\bigr)$.
And that is what the natural language explanation above says.}

So \citetalias{hott-book} never unnecessarily uses propositional truncation.
This gives its logic a different feel from the logic taught in philosophy
classes.
Our phrase "mathematically, there is no point to the forgetting" is called
using \emph{proof-relevant mathematics} (or proof-relevant logic) by
\citetalias{hott-book} (Section~1.11).

\subsection{The Classical Axiom of Choice}

Now we can also clarify the difference between the type theory axiom
of choice mentioned above and the classical axiom of choice.  The latter is
\citepalias[Lemma~3.8.2]{hott-book}, if $A$ is a set and $B(x)$ is a set
for all $x : A$, then
\begin{equation} \label{eq:set-theoretic-axiom-of-choice}
   \left( \prod_{x : A} \bigproptrunc{B(x)} \right) \to 
   \bigproptrunc{\prod_{x : A} B(x)}
\end{equation}
which is read just like the classical axiom of choice: if each $B(x)$
is inhabited, then so is $\prod_{x : A} B(x)$.
But we repeat that without the propositional truncation this is just true
(indeed, here, trivially true, because when we remove the vertical bars
both sides of the arrow are the same, and any function type $A \to A$ is
inhabited by the identity function).

\subsection{Homotopy Equivalence}
\label{sec:homotopy-equivalence}

The point of this section is that in HoTT isomorphism is too strong a concept.

\subsubsection{Homotopies of Functions}

In classical algebraic topology we have homotopies defined as continuous
functions $I \to X \to Y$ and $I \to I \to X \to Y$ and so forth, where
$I$ is the unit interval of the real line.

In HoTT we also use the notion of homotopies to talk about identity types,
but this is mere analogy.  A homotopy (analogically) is a path between paths,
a path between paths between paths and so forth, where path is just another
word (analogy) for a term of an identity type.

But \citepalias[Section~2.4]{hott-book} in HoTT we also use a defined notion
of the homotopy of two functions, which is a type.  Given
two dependent functions of the same type
$$
   f, g : \prod_{x : A} B(x)
$$
a \emph{homotopy} from $f$ to $g$ is a term of the type
\begin{equation} \label{eq:homotopy-type}
   (f \sim g) :\equiv \prod_{x : A} \left(f(x) =_{B(x)} g(x)\right)
\end{equation}
that is, the homotopy is a function that maps $x : A$ to a path
$p(x) : \left(f(x) =_{B(x)} g(x)\right)$.  The idea is that in HoTT
all functions are continuous (again, this is mere analogy, not defined
in terms of anything) so $p(x)$ is continuous in $x$ and also in movement
along the path.  So it is analogous to a classical homotopy.

The HoTT book cautions that \eqref{eq:homotopy-type} is not the identity
type $(f = g)$, so not a path
(rather a type of functions from points to paths).

\subsubsection{Quasi-Inverses}

Using the just-defined notion of homotopy, \citetalias[Definition~2.4.6
and~Section~4.1]{hott-book} defines for $f : A \to B$
\begin{equation} \label{eq:quasi-inverse-type}
   \qinv(f) :\equiv \sum_{g : B \to A}
   (f \circ g \sim \idfun_B) \times (g \circ f \sim \idfun_A)
\end{equation}
to be the type whose terms are \emph{quasi-inverses}.  The ``quasi''
is there to remind us that we have only homotopy rather than equality.
(If we replace the homotopies by equalities, then we get the usual
notion of isomorphism.  Or, if we work in homotopy category, this
is the notion of isomorphism for that category.)

Since a dependent sum type is an ordered pair and a product type is also
a pair, a term of $\qinv(f)$ has the form $(g, (G, H))$
where $G : (f \circ g \sim \idfun_B)$
and $H : (g \circ f \sim \idfun_A)$.
So this is the usual proof-relevant logic:
$g$ is a function and $(G, H)$, a pair of homotopies, is the proof
that $g$ is a quasi-inverse of $f$.

The trouble is, that (surprisingly) \eqref{eq:quasi-inverse-type}
is not a mere proposition
\citepalias[Theorem~4.1.3]{hott-book}.  So we could use propositional
truncation.  But (again perhaps surprisingly) the HoTT book does not.

\subsubsection{Bi-Inverses}

A function $f : A \to B$ is \emph{bi-invertible}
\citepalias[Definitions~4.2.7 and~4.3.1]{hott-book} if it has type
\begin{equation} \label{eq:bi-invertible-type}
   \binv(f) :\equiv \left( \sum_{g : B \to A} (f \circ g \sim \idfun_B)\right)
   \times \left( \sum_{h : B \to A} (h \circ f \sim \idfun_A)\right)
\end{equation}
So a term of this type is a pair of pairs $((g, G), (h, H))$
where $g$ and $h$ are functions and $G$ and $H$ are homotopies,
and $G$ proves that $g$ is a right inverse of $f$ and $H$ proves that
$h$ is a left inverse of $f$ (in homotopy category).

Unlike \eqref{eq:quasi-inverse-type}, \eqref{eq:bi-invertible-type} is
a mere proposition \citepalias[Theorem~4.3.2]{hott-book}.

\subsubsection{Homotopy Equivalences}

In its Sections~4.2 and~4.4 \citetalias{hott-book}
defines two other types that turn out to be (homotopy)
equivalent to bi-equivalence \eqref{eq:bi-invertible-type}.
Any could be used to define
\emph{homotopy equivalence}.\footnote{Exercise~3.8 {in} \citetalias{hott-book}
says $\isEquiv(f)$ could also be defined as $\proptrunc{\qinv(f)}$.}
The HoTT book chooses a different one from what we do.
We define $\isEquiv(f)$ to be the same as $\binv(f)$.

In the Notes for its Chapter~4 \citetalias{hott-book} explains that
the quasi-inverse property being different from homotopy
equivalence was well known in homotopy theory from its inception, long
before HoTT arose.  (But of course HoTT is a synthetic theory and classical
homotopy theory an analytic theory that needs real analysis and point set
topology to get off the ground.)

Exercise~4.6 {in} the \citetalias{hott-book} explains why
bi-invertibility is a good concept and cannot be replaced
by quasi-invertibility.

It is important to understand that $\isEquiv(f)$ is inhabited if and
only if $\qinv(f)$ is inhabited.\footnote{This is stated in the text
of Section~2.4 of \citetalias{hott-book} in items (i) and (ii) of the
list preceding equation (2.4.10), and proved in the text following that
equation.}  Thus the two concepts are logically equivalent in one sense,
but not in another because one is a mere proposition and the other is not.
Moreover, if $((g, G), (h, H)) : \binv(f)$, then both $g$ and $h$
are quasi-inverses of $f$ (one just needs different homotopies to prove that).


\subsection{Universes and Univalence}
\label{sec:univalence}

In HoTT we allow types to have types.  This again raises the specter of
Russell's paradox.  So we collect types into universe types $\universe$,
which are closed under the operations of type theory
(those in Table~\ref{tab:foo}).
Then we can consider a dependent type as a function $B : A \to \universe$.

To be careful we cannot have $\universe : \universe$ because that
would allow Russell's paradox.  So we might have to have a sequence of
ever larger universes.  (as Russell and Whitehead did).
But normally, we do not need to worry about that
\citepalias[Section~1.3]{hott-book} and won't here.

Then for $A, B : \universe$ we can consider the equality
type $(A =_\universe B)$.  And it is provable in MLTT
\citepalias[Lemma~2.10.1]{hott-book} for $A, B : \universe$ there
exists a function defined in the proof of the cited lemma
$$
   \text{idtoequiv} : (A =_\universe B) \to (A \simeq B)
$$
Voevodsky proved that it is consistent\footnote{Except, of course,
G\"{o}del's inconsistency theorem says that any theory strong enough
to contain arithmetic cannot prove its own consistency, so what
Voevodsky constructed was a model of HoTT with univalence in set theory.
Hence these theories are equiconsistent.  \citet*{bezem-coquand-huber}
constructed a model of HoTT with univalence in a constructive type theory.
So Voevodsky used the axiom of choice, but \citet{bezem-coquand-huber}
did not, nor even LEM.}
to assume as an axiom that this function is an equivalence
$$
   (A =_\universe B) \simeq (A \simeq B)
$$
which says $\text{idtoequiv}$ has a quasi-inverse
$$
   \text{ua} : (A \simeq B) \to (A =_\universe B)
$$
($\text{ua}$ for univalence axiom).  And this function can be used to
transfer properties from $A$ to $B$.\footnote{The rules for $\text{ua}$,
its elimination rule, propositional computation rule, and propositional
uniqueness principle, discussed in \citetalias[Section~2.10]{hott-book}
do this.}

So this is the high water mark of the HoTT view of equivalence.  It
is sometimes misunderstood as saying that isomorphic types are equal,
but that is wrong in two ways.  First, isomorphism is a stronger property
than (homotopy) equivalence.  Second, it is not that equivalence is being
reduced to equality (some sort
of skeletality condition) but rather that equality (in the universe)
is being expanded to be (homotopy) equivalent to (homotopy) equivalence.

One should not think of the univalence axiom as some sort of broad general
principle to be used everywhere in philosophy but rather as a technical result
needed to make identity types in MLTT work the way we want them to in HoTT.

The univalence principle when applied to mere propositions
(propositional extensionality) was present in the early type theories
of Russell and Church \citepalias[Notes for Chapter~2]{hott-book}.

\subsection{Law of Excluded Middle}

We can now also clarify the status of LEM in
HoTT.  The following statement
$$
   \prod_{A : \universe} \bigl(A + (A \to 0)\bigr)
$$
is incompatible with the univalence axiom
\citepalias[Theorem~3.2.2]{hott-book}.  So we cannot have that.

But we can consistently assume \citepalias[Section~3.4]{hott-book}
\begin{equation} \label{eq:propositional-lem}
   \prod_{A : \universe} \Bigl(\isProp(A) \to \bigl(A + (A \to 0)\bigr)
   \Bigr)
\end{equation}
However, we do not routinely assume this (sometimes denoted $\text{LEM}_{-1}$)
because it makes our theorems not applicable to large parts of contemporary
mathematics.

As a simple example, notice that this document always says inhabited rather
than nonempty, and
equivalence of these two concepts is literally double negation removal,
which is equivalent to LEM.  Often, just changing the way you say things
allows for constructive mathematics.

\subsection{Function Extensionality and Univalence}

The following is provable in MLTT. Given a type $A$,
a type family $B : A \to \universe$,
and functions $f, g : \prod_{x : A} B(x)$,
\begin{equation} \label{eq:happly}
   (f = g) \to \prod_{x : A} \bigl( f(x) =_{B(x)} g(x) \bigr)
\end{equation}
that is, $f = g$ implies $f(x) = g(x)$ for all $x$.
The reverse implication is not provable in MLTT.

The assumption that the inhabitant of \eqref{eq:happly} that can be proved
to exist is a homotopy equivalence
is called the axiom of \emph{function extensionality}
\citepalias[Axiom~2.9.3]{hott-book}.
It is hard to do any mathematics without this axiom, which is much more
fundamental than LEM or "true" equality (equality in the logic).
All of mathematics assumes it.
However, HoTT does not assume this as an axiom because it implied
by the univalence axiom \citepalias[Section~4.9]{hott-book}.

Of course, first-order classical logic plus set theory also
proves the function extensionality property because it just defines
functions to be extensional relations.

\subsection{Homotopy Levels}
\label{sec:homotopy-levels}

There is a pattern in the definitions of (mere) propositions
and (mere) sets in HoTT.

Start with contractibility.  A type $A$ is contractible
\citepalias[Section~3.11]{hott-book}
if it has exactly one path component, that is,
\begin{equation} \label{eq:contractible}
   \isContr(A) :\equiv \sum_{a : A} \prod_{b : A} (a =_A b)
\end{equation}
read there exists $a$ in $A$ such that for all $b$ in $A$ we have $a =_A b$.

Now for the pattern.  We define homotopy-$n$-types
\citepalias[Chapter~7]{hott-book}
by
\begin{align}
   \text{is-$(- 2)$-type}(A) & :\equiv \isContr(A)
   \\
   \text{is-$(n + 1)$-type}(A)
   & :\equiv
   \prod_{x, y : A} \text{is-$n$-type}(x =_A y),
   \qquad n \ge -2
   \label{eq:type-recursion}
\end{align}
The bizarre idea to start counting at $- 2$ is to match up with how
homotopy levels were defined in homotopy theory before HoTT was invented.

It is then a theorem \citepalias[Lemma~3.11.10]{hott-book} that
$$
   \isProp(A) \simeq \text{is-$(-1)$-type}(A)
$$
Then it is obvious that
$$
   \isSet(A) \simeq \text{is-$0$-type}(A)
$$
And we can go on naming levels
$$
   \text{isGroupoid}(A) :\equiv \text{is-$1$-type}(A)
$$
This makes sense because we think of $a : A$ as the objects of the groupoid
and $p : (x =_A y)$ as morphisms, in which case the laws for identity types
(Section~\ref{sec:identity} above) are the category laws
(Section~\ref{sec:category} above) with the additional fact that any such $p$
is invertible (hence an isomorphism when considered a morphism) by the symmetry
requirement on paths.

But then we run out of names.  A 2-type is (if you have to make up
a name) a 2-groupoid, (having nontrivial paths between paths).
A 3-type is (if you have to make up
a name) a 3-groupoid, (having nontrivial paths between paths between paths).
And so on.

MLTT cannot prove whether or not any homotopy $n$-levels above sets exist.
HoTT can prove they do exist
using univalence \citepalias[Example~3.1.9]{hott-book}.
Not every type is an $n$-type for some $n$.
An example is the two-sphere $S^2$ (like the surface of a globe).
Such a type is then called an $\infty$-groupoid in (higher)
category theory.\footnote{The HoTT book does not actually prove this about
$S^2$, but does construct \citepalias[Example~8.8.6]{hott-book} a type that
is not a homotopy $n$-type for any $n$.}

HoTT has \citepalias[Section~7.3]{hott-book} an operation
$A \mapsto \proptrunc{A}_n$ that gives for any type $A$ its best approximation
(in some sense) that is a homotopy $n$-type.  This construction is due
to Voevodsky and corresponds to the $n$-th Postnikov section in classical
homotopy theory.  When applied to a type $A$ that is not a homotopy $n$-type
for any $n$, this gives $\proptrunc{A}_n$ inhabiting homotopy level $n$.
Propositional truncation (Section~\ref{sec:truncation} above) is the
same as $\proptrunc{\fatdot}_0$.

All of this requires univalence (or something like it).
It is consistent with MLTT to assume that all types are sets
\citepalias[Remark~8.1.17 and Section~7.2]{hott-book}.
An equivalent assumption is that for $x, y : A$ and $p, q : x =_A y$
we have $p = q$, which is called \emph{uniqueness of identity proofs}
(UIP).\footnote{Propositions-as-types says $p$ and $q$ are proofs
of the proposition $x = y$.}
Another equivalent assumption is Streicher’s Axiom K, which asserts
for $x : A$ and $p : x =_A x$ we have $p = \refl_x$.
We stress that HoTT assumes none of UIP or Axiom K or all types are sets.
The whole point of HoTT is not to dumb down MLTT to only do mathematics
as it was understood in 1900.

\subsection{The Natural Numbers and Finiteness}

In HoTT the natural number type $\nats$ has two constructors
\begin{align*}
   0_\nats & : \nats
   \\
   \successor & : \nats \to \nats
\end{align*}
The first subscripts the term $0_\nats$ to distinguish it from the empty
type 0.  Every term belongs to a distinct type.
The second is a function.  This makes $\nats$ to be the sequence
$$
   0_\nats,
   \successor(0_\nats),
   \successor(\successor(0_\nats)),
   \successor(\successor(\successor(0_\nats))),
   \ldots
$$
The idea is that an inductive type like this is freely generated by
the constructors.  But to actually use the type we also need elimination
rules.  For an inductive type like this, the most general elimination rule,
called the \emph{dependent eliminator} or the \emph{induction principle},
is a function
\citepalias[Section~1.9]{hott-book}
\begin{equation}
\label{eq:natural-numbers-induction}
   \inductor_\nats : \prod_{C : \nats \to \universe}
   C(0_\nats) \to \left( \prod_{n : \nats} C(n) \to C(\successor(n)) \right)
   \to \prod_{n : \nats} C(n)
\end{equation}
having defining equations
\begin{subequations}
\begin{align}
   \inductor_\nats(C, c_0, c_s, 0_\nats) & :\equiv c_0
   \label{eq:nats-induction-base}
   \\
   \inductor_\nats(C, c_0, c_s, \successor(n)) & :\equiv
   c_s\bigl(n, \inductor_\nats(C, c_0, c_s, n) \bigr)
   \label{eq:nats-induction-succ}
\end{align}
\end{subequations}
This is a little hard to decode, so let's be pedantic.
\begin{itemize}
\item $\inductor_\nats$ is a $\Pi$-type,
    hence a dependent function whose argument is the subscript
    of the $\Pi$, here $C : \nats \to \universe$ and the value is a function
$$
   C(0_\nats) \to \left( \prod_{n : \nats} C(n) \to C(\successor(n)) \right)
   \to \prod_{n : \nats} C(n)
$$
   (the rest of the expression).
   So $\inductor_\nats(C)$ is a term of this type.
\item And this is a function whose argument is $c_0 : C(0_\nats)$ and
    whose value is a function
$$
   \left( \prod_{n : \nats} C(n) \to C(\successor(n)) \right)
   \to \prod_{n : \nats} C(n)
$$
    So $\inductor_\nats(C)(c_0)$ is a term of this type.
\item And this is a function $c_s$ whose argument is a $\Pi$-type
    (another dependent function) which for $n : \nats$ is a function
    $c_s(n) : C(n) \to C(\successor(n))$ and whose value is yet another
    $\Pi$-type.  So $\inductor_\nats(C)(c_0)(c_s)$ is a term of type
    $n \mapsto C(n)$.
\item
    So $\inductor_\nats(C)(c_0)(c_s)(n)$ is a term of type $C(n)$.
\end{itemize}
Hence, in summary, we have a function $\inductor_\nats(C)(c_0)(c_s)(n)$,
which we can also
write as a function of four arguments
$\inductor_\nats(C, c_0, c_s, n)$\footnote{Currying
(footnote~\ref{foot:currying} above).} and which maps to $C(n)$.
The defining equations \eqref{eq:nats-induction-base} and
\eqref{eq:nats-induction-succ} say that if we have
\begin{align*}
   c_0 & : C(0)
   \\
   c_s & : \prod_{n : \nats} \bigl( C(n) \to C(\successor(n)) \bigr)
\end{align*}
Then we have $\prod_{n : \nats} C(n)$, that is, the function
$\inductor_\nats(C, c_0, c_s) : n \mapsto C(n)$
given by the induction principle.
Reading the $\Pi$-types as logical for-all, we have classical mathematical
induction.

Now what are we to make of this?  This is exactly the same as
Peano's axioms for arithmetic\footnote{Actually, Peano needed two other
axioms saying that $0_\nats \neq \successor(n)$ for any $n$
and $\successor$ is an injective function.
But these are provable in HoTT
\citepalias[Section~2.13]{hott-book}, so we do not add them to the rules
for type $\nats$.}
except for here $C : \nats \to \universe$
is any type family, whereas the Peano axioms only apply to mere propositions,
that is, $C(n)$ can here be any type, but for the Peano axioms it must be
a mere proposition.  There is also the difference that the Peano axioms
are based on classical logic and this is based on propositions-as-types
(intuitionistic, proof-relevant logic).

The point is that HoTT proves all and only what Peano arithmetic proves
about the natural numbers.  And all of this is intensional and synthetic
rather than extensional and analytic.  There is nothing in any of this
about an intended model or any model or models.  Just constructions/proofs.

Another way to say this is that the induction principle says that all you
have to do to prove $\prod_{n : \nats} P(n)$ is to verify $P(0)$ and
and $P(n) \to P(\successor(n))$ for each $n : \nats$.
There are no other cases to check.
And how do we know that?  By the rules for type $\nats$.

\subsection{Path Induction}
\label{sec:path-induction}

Identity types have an induction principle
described by \citetalias[Section~1.12]{hott-book} as follows.
Given a dependent type family
$$
   C : \prod_{x, y : A} (x =_A y) \to \universe
$$
and a dependent function
$$
   c : \prod_{x : A} C(x, x, \refl_x)
$$
there is a dependent function
$$
   f : \prod_{x, y : A} \prod_{p : (x =_A y)} C(x, y, p)
$$
such that
$$
   f(x, x, \refl_x) :\equiv c(x)
$$
Put in the form of a single function, the induction principle is
$$
   \inductor_{=_A} :
   \prod_{C : \prod_{x, y : A} (x =_A y) \to \universe}
   \left( \prod_{x : A} C(x, x, \refl_x) \right)
   \to
   \prod_{x, y : A} \prod_{p : (x =_A y)} C(x, y, p)
$$

Admittedly, this is hard to read and hard to follow
(the \citetalias{hott-book} says ``most subtle'').
Read sloppily, path induction says that if you prove a property for
reflexivity paths, then you have proved it for all paths.
And (also confusingly) many proofs read this way: by induction we may
assume $p$ is $\refl_x$, then whatever we are supposed to check checks,
QED.

The subtlety \citepalias[Remark 1.12.1]{hott-book} is that
the input is the dependent family of types $(x, y) \mapsto (x =_A y)$,
so $x$ and $y$ must vary freely over $A$.
If $x$ or $y$ is fixed (not varying) then checking something about $\refl_x$
proves nothing.

Another way to say the same thing is that proofs by path induction start
with an arbitrary path $p : (x =_A y)$ and prove that if you can do something
to $p$ when $p \equiv \refl_x$, then you can do the same thing when $p$ is
arbitrary.  But you have to start with having the arbitrary $p$.

There is a variant of path induction called \emph{based path induction}
in which only one endpoint of the path is freely varying while the other
end is fixed; \citetalias{hott-book} then proves
this is equivalent to general path induction.  But fixing both endpoints
still proves nothing.

Chapter~2 of \citetalias{hott-book} shows that all of the properties
of identity types discussed in Section~\ref{sec:identity} above are
provable from path induction.  And many more properties besides.

For example, Section~2.2 of \citetalias{hott-book} is titled ``functions
are functors'' although, strictly speaking, they are not.  What is meant
is that functions act functorially on paths.
Lemma~2.2.1 of \citetalias{hott-book} says that for any function $f : A \to B$
and any $x, y : A$ there is a function
\begin{subequations}
\begin{equation} \label{eq:path-induction}
   \text{ap}_f : (x =_A y) \to (f(x) =_B f(y))
\end{equation}
such that for each $x : A$ we have
\begin{equation} \label{eq:path-induction-too}
   \text{ap}_f(\refl_x) \equiv \refl_{f(x)}
\end{equation}
\end{subequations}

Thus every path $p$ from $x$ to $y$ is lifted to a path $\text{ap}_f(p)$
from $f(x)$ to $f(y)$.  And all of this is synthetic rather than analytic,
just there because of path induction.  We can read $\text{ap}_f$ as meaning
applying $f$ to a path.  Of course, the argument of $f$ is not a path;
$f$ takes arguments $a : A$, and paths are not made out of points of $A$.
Their endpoints are points of $A$, but that is all.\footnote{But if a path were
(as it was before HoTT) a function $p : I \to A$, where $I$ is the unit interval
of the real line, then this lifted path would just be $f \circ p$.  Moreover,
the HoTT book and others also use $f(p)$ instead of $\text{ap}_f(p)$ when
they think this abuse of notation will cause no confusion.}

As an example of the opaqueness of path induction, we copy
\citetalias{hott-book}'s second proof of the application-to-path lemma.
To define $\text{ap}_f$ for all $p : x = y$, it suffices, by induction,
to assume $p$ is $\refl_x$, in which case $\text{ap}_f(p)$ is
given by \eqref{eq:path-induction-too}.
This is a valid application of path induction because $x$ and $y$ are
free rather than fixed.

It is remarkable that \citet{martin-lof-mltt} figured out path induction
from first principles of constructive logic without any recognition of
the homotopy motivation.  \citet{ladyman-presnell} are not philosophically
satisfied with \citeauthor{martin-lof-mltt}'s justification and provide
their own pre-mathematical (and thus philosophical)
justification.\footnote{Their argument is a combination of what
\citetalias[Section~1.12]{hott-book} calls ``indiscernability of identicals''
(equals may be substituted for equals) and (Remark~1.12.1) ``a path
$p : x = x$ may be not equal to reflexivity as an element of $( x = x )$,
but the pair $(x, p)$ will nevertheless be equal to the pair $(x, \refl_x)$
as elements of $\sum_{(y:A)} (x = y)$.  And this is just the point we said
about at least one endpoint of the path must vary freely to use path induction.
\citeauthor{ladyman-presnell} recast this so as to satisfy them as being
pre-mathematical, hence fundamental.  It is not clear what pre-mathematical
is supposed to mean given that mathematics is as old as philosophy and a major
motivation for it, as the enscription
over the entrance to Plato's Academy indicated.
Frege, Russel, and other creators of analytic philosophy also knew so much
mathematics that it is unclear what pre-mathematical intuitions they may
have had.}
The fact that identity types in MLTT can be as complicated as this
section describes was first shown by \citet{hofmann-streicher}.

\subsection{Higher Inductive Types}
\label{sec:higher-induction}

More complicated instances of induction are called \emph{higher induction}
when there are constructors for not only terms of a type but also paths
(identitifications).

We have already met one of these, propositional truncation,
which has two constructors (repeated from Section~\ref{sec:truncation} above)
\begin{align*}
   \termtrunc{\fatdot} & : A \to \proptrunc{A}
   \\
   p & : \text{$(x =_{\proptrunc{A}} y)$ whenever $x, y : \proptrunc{A}$}
\end{align*}
The first constructor says\
\begin{itemize}
\item if $A$ is a type, then $\proptrunc{A}$ is a type, and
\item if $a : A$, then $\termtrunc{a} : \proptrunc{A}$.
\end{itemize}
The second constructor just formalizes what we already said
(Section~\ref{sec:truncation} above):
we do propositional truncation by redefining
equality to make all points in the truncation equal.

Now for something new and nontrivial, some synthetic algebraic topology.
The (topological) circle, denoted $S^1$ because it is the one-dimensional
analog of the (two-dimensional surface of a three-dimensional) sphere
(denoted $S^2$) is defined in HoTT as the type with two constructors
\begin{align*}
   \sbase & : S^1
   \\
   \sloop & : \sbase =_{S^1} \sbase
\end{align*}
This is remarkable because in analytic geometry, you need the real numbers
$\real$ and the plane $\real^2$ to define a circle.  And even in synthetic
(Euclidean) geometry, you need some set-up (Euclid's axioms).
Here we have only types and identity types.

One might wonder how we know that $\sloop$ is not $\refl_\sbase$,
but this is provable constructively using univalence
\citepalias[Lemma~6.4.1]{hott-book}.

Now by equality being reflexive, symmetric, and transitive, we can
inductively define
\begin{alignat*}{2}
   \sloop^0 & :\equiv \refl_\sbase
   \\
   \sloop^1 & :\equiv \sloop
   \\
   \sloop^{n + 1} & :\equiv \sloop^n \cdot \sloop,
        & \qquad & n \ge 1
   \\
   \sloop^{- (n + 1)} & :\equiv \sloop^{- n} \cdot \sloop^{-1},
        & \qquad & n \ge 1
\end{alignat*}
Let $\ints$ denote the type of integers \citepalias[which we haven't defined
but is defined in the usual way from the
natural numbers,][Section~6.10]{hott-book}.  Then we have a map
$\ints \to \pi_1(S^1)$ given by $n \mapsto \sloop^n$, where $\pi_1(S^1)$
denotes the fundamental group of $S^1$ whose elements are loops with the
same basepoint and whose group operation is loop concatenation.

We know from algebraic topology that the fundamental group of $S^1$ is $\ints$
(up to isomorphism of groups, with the group operation for $\ints$ being
addition) \citepalias[Table~8.1]{hott-book}.  Thus we want to prove that
the map $n \mapsto \sloop^n$ is a group isomorphism.
For the first step, it is easily verified that it is a group homomorphism,
that $\sloop^m \cdot \sloop^n = \sloop^{m + n}$.

And then comes the shocker.  Or perhaps not shocking at all.  Math is hard.
In order to prove this function is a homotopy equivalence, hence
and equality by univalence,
\citetalias[Section~8.1]{hott-book} takes more than seven pages of very dense
mathematics (they do give three proofs though, one the classical
proof from algebraic topology translated into HoTT,
and two HoTT-friendly proofs).
They also prove that the loop space of $S^1$ is a set, so all higher homotopy
groups (above the first discussed above) are trivial.
And everything they do is synthetic homotopy theory.
There is no use whatsoever of real analysis or point set topology.

For our next trick, we define $S^2$ as a higher inductive type
with constructors
\begin{align*}
   \sbase & : S^2
   \\
   \ssurf & : \refl_\sbase =_{\left(\sbase =_{S^2} \sbase\right)} \refl_\sbase
\end{align*}
That is, like $S^1$ we have a basepoint called $\sbase$.
But, unlike $S^1$, we skip one-dimensional paths and put in a two-dimensional
surface called $\ssurf$, which is a path between paths ($\refl_\sbase$ to
itself).  That's it.  Like $S^1$ there are no more constructors or rules.

But, unlike $S^1$, this type is much more complicated homotopically.
It has nontrivial homotopy groups of all orders
\citepalias[Table~1]{hott-book}.
And, of course, this means it has nonreflexive identities (paths) of all
orders.  So where do they come from if there are no constructors for them?

And, on the same theme, we can ask the opposite question about $S^1$.
Why does $S^1$ have only one nontrival homotopy group?  Neither question
has an obvious answer, but \citetalias{hott-book} does give a complete
answer to the question about $S^1$.  The answer comes from classical
algebraic topology using the fact that the real number system considered
as a topological space is the universal cover of $S^1$
\citepalias[Section~8.1.3]{hott-book}.  But at this point in
\citetalias{hott-book} the real number system has not been constructed;
that is its Chapter~11.  So here they use the \emph{homotopical reals}
\citepalias[Section~8.1.5]{hott-book}, a higher inductive type $\real_h$ with
constructors
\begin{align*}
   c & : \ints \to \real_h
   \\
   p(z) & :
   \text{$c(z) =_{\real_h} c(\successor(z))$ whenever $z : \ints$}
\end{align*}
that is, we start with $\ints$ (which has no nonreflexive paths in its
path spaces)
and add paths between predecessors and successors
(which implies paths between each pair of points).
This construction is useless for real and functional analysis,
but great for topology.

The answer to the question about $S^2$ is much more complicated.
It has to do with the fact that paths between paths between paths
and so forth \emph{ad infinitum} do not just exist, but have complicated
interrelationships.  Chapters~2 and~8 of \citetalias{hott-book} are
full of theorems about that, and there is a lot more they could have
proved but left out \citepalias[Remark~2.1.5]{hott-book}.
So it should not be a surprise that the one constructor that puts
a two-dimensional path in $S^2$ implies the existence of paths of all higher
dimensions.  The surprise is that $S^1$ does not have any such higher
dimensional paths that are not equal to $\refl$.
Sometimes mathematics is just like that: very complicated with no simple
story.  Prime numbers have to be discovered one at a time.  There is no
formula.  Homotopy groups of $S^2$ have to be discovered one at a time.
There is no formula \citepalias[discussion of Table~8.1]{hott-book}.

\subsection{Inductive-Inductive Types}

The HoTT book closes with (Chapter~11) the construction of the real number
system and Conway's surreal number system.  It turns out that one cannot
constructively prove that the Cauchy real numbers $\real_c$ and the Dedekind
real numbers $\real_d$ are equivalent (it can be proved that $\real_c$ is
a subset of $\real_d$).
One can prove $\real_c \simeq \real_d$ if LEM or the axiom of countable choice
are assumed \citepalias[Corollary~11.4.3]{hott-book}.
\citet{bishop} did accept the axiom of countable choice, so he had only one
real number system.

But \citet{bishop} did not accept LEM, so there were still differences between
his real number system and the one from mainstream real analysis (which does
accept LEM).

A type has \emph{decidable equality} if
LEM holds for propositions of the form $a = b$ for terms $a$ and $b$ of the
type \citepalias[Definition~3.4.3(iii)]{hott-book}.
The natural numbers, integers, and rational numbers
have decidable equality \citepalias[Section~11.1]{hott-book}.
The real numbers do not.  For example $\neg (x = 0)$ does not imply $x$ has
a multiplicative inverse $1 / x$.  Instead a stronger condition
\emph{apartness}, defined as $\lvert x \rvert > 0$, is required
\citepalias[Section~11.2.1]{hott-book}.
\citet[Section~3.2]{bauer-constructive} proves that no computer program
can verify whether or not an arbitrary real number is zero
(to an infinite number of decimal places).
Hence no computer program can verify LEM.

This is related to the computer programming principle that one should never
compare the computer's so-called real numbers (which are, of course, only
approximations) for equality.  Different computations that should come out
equal if real real numbers were used do not due to inexactness of computer
arithmetic.  One can show by formal proof that two formulas or two computer
programs should correspond to or compute equal results (up to whatever
degree of approximation they are evaluated to).  But one cannot rely on
computed results (representations of numbers) being equal.  Thus intuitionistic
logic corresponds to the way mathematics is actually used in computing.

One might wonder whether this makes math too hard to do.
It was certainly thought so before \citet{bishop},
but he showed a tremendous amount of real analysis, functional analysis,
and measure theory could be done without LEM.  One just has to be a bit
more careful, and this is the same care that is required for
accurate computations.

\citet[Section~5.1]{bauer-constructive} discusses the intermediate value
theorem from calculus, which classically (with LEM) says if $x < y$
and $f(x) \le 0 \le f(y)$ and $f$ is continuous, then $f(z) = 0$ for
some $z$ such that $x \le z \le y$.
This cannot be proved without LEM, and the reason is what we said above,
you cannot tell (constructively) whether or not a real number is zero
(and are a fool to even try).

But there are useful modifications of the classical theorem that
are constructively provable.  One uses approximation
\citep[Theorem~5.4]{bauer-constructive}: if $x < y$ and $f(x) \le 0 \le f(y)$
and $f$ is continuous and $\varepsilon > 0$,
then  $\lvert f(z) \rvert < \varepsilon$ for some $z$ such that $x \le z \le y$.

Returning to HoTT, the Dedekind reals are constructed more-or-less as in
classical mathematics.  Terms of the type are Dedekind cuts
(some care is required in definitions to avoid using LEM).

The Cauchy reals are more complicated and more interesting.
No method of constructing the Cauchy reals was known before that of
\citetalias[Section~11.2,]{hott-book} that did not use LEM or the axiom
of countable choice.  The HoTT book uses an inductive-inductive definition,
which simultaneously defines the type $\real_c$ and a ternary relation
$\rats_+ \to \real_c \to \real_c \to \text{Prop}$, where $\rats$ denotes
the rational number system and $\rats_+$ its subset of positive elements,
and $\sim(\varepsilon, x, y)$ is usually written $x \sim_\varepsilon y$.

The idea of the Cauchy reals is that they are the Cauchy completion of
the rationals, the set of limits of Cauchy sequences of rationals.
This is a special case of completion of metric spaces.
A sequence $n \to x_n$ is Cauchy if $\lvert x_m - x_n \rvert \to 0$ as
$\min(m, n) \to \infty$.
A metric space is \emph{complete} if every Cauchy sequence converges.
Classically (using the axiom of choice) one constructs the completion
of an arbitrary metric space as the set of equivalence classes of Cauchy
sequences that converge to the same point (where same means distance zero).
There are a few problems in applying this idea to the construction of the
real numbers.  First, a distance is a real-valued function, so we cannot
define that before we define the real numbers.  But we can apply
a rational-valued distance.  That is what $\rats_+$ is doing in the
the definition of $x \sim_\varepsilon y$.
Then \citepalias[Definition~11.3.2]{hott-book}
the Cauchy number system $\real_c$ has constructors
\begin{itemize}
\item \textbf{rational points:} if $q : \rats$,
    then there is $\text{rat}(q) : \real_c$.
\item \textbf{limit points:} if $x : \rats_+ \to \real_c$
    satisfies $\prod_{\delta, \varepsilon : \rats_+}
    (x_\delta \sim_{\delta + \varepsilon} x_\varepsilon)$
    (such an $x$ is called a \emph{Cauchy approximation}),
    then there exists $\text{lim}(x) : \real_c$.
\item \textbf{paths:} If $u, v : \real_c$
    such that $\prod_{\varepsilon : \rats_+} (u \sim_\varepsilon v)$,
    then there exists a path of type $(u =_{\real_c} v)$.
\end{itemize}
and the relation $\sim$ has constructors, where in all constructors
$q, r : \rats$ and $\delta, \varepsilon, \eta : \rats_+$ and $u, v : \real_c$
and $x$ and $y$ are Cauchy approximations (as defined above)
\begin{itemize}
\item If $- \varepsilon < q - r < \varepsilon$,
    then $\text{rat}(q) \sim_\varepsilon \text{rat}(r)$.
\item If $\text{rat}(q) \sim_{\varepsilon - \delta} y_\delta$,
    then $\text{rat}(q) \sim_\varepsilon \text{lim}(y)$.
\item If $x_\delta \sim_{\varepsilon - \delta} \text{rat}(r)$,
    then $\text{lim}(x) \sim_\varepsilon \text{rat}(r)$.
\item If $x_\delta \sim_{\varepsilon - \delta - \eta} y_\eta$,
    then $\text{lim}(x) \sim_\varepsilon \text{lim}(y)$.
\item Any type $(u \sim_\varepsilon v)$ is a mere proposition
    by propositional truncation.
\end{itemize}
(It is assumed when $\sim_{\varepsilon - \delta}$ and $\sim_{\varepsilon
- \delta - \eta}$ appear that the subscripts are positive (as is required
by the definition of $\sim$).)

With all of that, the definition is entirely constructive
with no use of LEM or any choice axioms.

The notes for Chapters~6 and~11 {in} \citetalias{hott-book} say the theories
of higher inductive types and inductive-inductive types were not fully
worked out at the time that book was written and (AFAIK) not now either,
although work is ongoing.  They do provide good illustrative examples.

\section{The Smorgasbord}

We want to emphasize that in HoTT you pick and choose what you want.

\begin{itemize}
\item If you want LEM, you can assume that, \eqref{eq:propositional-lem} above.
\item If you want the axiom of choice, you can assume that,
    \eqref{eq:set-theoretic-axiom-of-choice} above.
\item If you want the countable axiom of choice, you can assume that,
    \eqref{eq:set-theoretic-axiom-of-choice} above with $A$ replaced by $\nats$.
\item If you want to assume every type is a mere set, also called UIP or
    Axiom K (Section~\ref{sec:homotopy-levels} above), you can assume that.
\item Conversely, if you want to assume that not every type is a mere set,
    then you assume the univalence axiom.
\end{itemize}

\citet[about 6 minutes into the 4th lecture]{bauer-hott} goes into a long
aside about ``structure versus property'' sometimes there is structure that
you don't want.  How do we get rid of it in HoTT?  His example is group
theory.  Groups are tuples $(G, m, i, e)$ where $G$ is the carrier of the
group (the set if you like), $m$ is the binary operation (multiplication or
addition), $i$ is the inverse operation, and $e$ is the identity element.
The way we write tuples in HoTT is iterated dependent pairs, so
\begin{equation*}
   \text{Group} :\equiv
   \sum_{G : \universe}
   \sum_{m : G \to G \to G}
   \sum_{i : G \to G}
   \sum_{e : G}
   \text{isGroup}(G, m, i, e)
\end{equation*}
And then one might try to define
\begin{multline*}
   \text{isGroup}(G, m, i, e) :\equiv
   \left( \prod_{x, y, z : G} m(m(x, y), z) = m(x, m(y, z)) \right)
   \\
   \times
   \left( \prod_{x : G} \bigl( m(x, i(x)) = e \bigr)
   \times \bigl( m(i(x), x) = e \bigr) \right)
   \\
   \times
   \left( \prod_{x : G} \bigl( m(x, e) = x \bigr)
   \times \bigl( m(e, x) = x \bigr) \right)
\end{multline*}
(the first big parens is the associative law, the second the
existence and properties of inverses, the third the
existence and properties of identities).
But this is wrong (\citealp{bauer-hott}
and \citetalias[Remark~6.11.2]{hott-book}).
First, it is a constructive or intuitionistic group because we have not
assumed LEM or the axiom of choice, but we won't worry about that.
Second, it is not a group as studied in group theory because it has (or at
least we have not ruled out) paths between paths, paths between paths between
paths, and so forth.  This $G$ is an $\infty$-group.
If we want a group, we need
\begin{multline*}
   \text{isGroup}(G, m, i, e) :\equiv
   \isSet(G)
   \times
   \left( \prod_{x, y, z : G} m(m(x, y), z) = m(x, m(y, z)) \right)
   \\
   \times
   \left( \prod_{x : G} \bigl( m(x, i(x)) = e \bigr)
   \times \bigl( m(i(x), x) = e \bigr) \right)
   \\
   \times
   \left( \prod_{x : G} \bigl( m(x, e) = x \bigr)
   \times \bigl( m(e, x) = x \bigr) \right)
\end{multline*}
The $\isSet(G)$ imposes the UIP property on this type.
We saw a similar move in the construction of the Cauchy reals, where
the last axiom for $\sim$ imposes the UIP property on $\real_c$.

So HoTT does not make classically trained mathematicians do anything they
don't want to.  It just makes them be explicit about it.
It makes a lot fewer default assumptions than they are used to.

\section{Set Theory as Spaghetti Code}

When computers were just invented, coding was an anything-goes affair.
There were no computer languages, just op (operation) codes which when
combined with addresses (numbers indicating memory locations)
made machine-language
instructions.  But everything, instructions and data were just bits (zeros
or ones).  What they represented (other than the op codes and memory addresses)
was completely up to the programmer.  There was no distinction between
instructions and data.  A program could rewrite its own code on the fly.

It did not take long before it was figured out that this is a horrible way
to program.  People could do it.  But it took intense concentration, and
did not scale.  Even genius programmers could not do it consistently, and
teams of ordinary programmers working on large projects often made such
a mess that the system was never delivered (this was the so-called
software crisis recognized in the 1960's).

Better ways to program were devised.  Firstly, programs were no longer allowed
to rewrite their own code.  Operating systems enforce the distinction.
This disallows a few useful techniques but makes programs much easier to
reason about.

Secondly, the ``structured programming'' paradigm, exemplified by many
high-level languages of the 1970's, such as Pascal, greatly restricted
control flow.  A famous (or notorious, depending on point of view) paper
\citep{go-to-harmful}
argued that the go-to statement should never be used, only better behaved
(and easier to reason about) structures like if-then-else and loops that
only have exit tests at the top or bottom (the difference is that if the
test is at the top, the loop body may be executed zero times).
The pejorative term that this paradigm had for previous practice was
``spaghetti code'' as tangled as a plate of spaghetti.

Thirdly (or perhaps thirdly and fourthly), type systems were added to most
computer languages.  This was a large part of both the ``object oriented''
and ``functional programming'' paradigms.  Most modern statically compiled
computer languages have relatively strong type systems.  Variables are
not allowed to refer to anything at all
but rather to objects of some specific type.
This makes it possible for the compiler to catch many errors.  If a program
does not type check, it is obviously wrong.

Thus types are not some weird thing only used in logic and type theory.
They are used everywhere in modern programming languages and are the
subject of large books \citep{pierce}.
But these subjects have actually merged in modern proof assistants.
One of the first strongly typed languages, ML, was designed to write proof
assistants in.  But its type system (Hindley–Milner) was taken over into
languages not for proof assistants, like Haskell, OCaml, and Scala.
Later languages for proof assistants, like Rocq, Agda, and Lean,
also use type theory.  It is possible to write a proof assistant based
on classical logic and set theory (Mizar), but that is not the trend.
Some proof assistants (HOL Light and Rocq) are written in OCaml and
can verify (some) programs written in OCaml.  Hence they can be used
to formally verify the correctness of working computer programs.
The proof assistant ACL2 is written in Common LISP and can formally
verify computer programs written in that language.

From all of this we can see that, in mathematics, the point of set theory
is to destroy the distinction between terms and types (both become sets)
or to destroy the idea of types at all.  Mathematicians have never informally
reasoned that way.  Euclid thought points and lines were different types of
things.  \citet{leinster} discusses how unnatural it is to ask the question:
What are the elements of $\sqrt{2}$?  If everything is a set, then this set
has elements.  But mathematicians never think that way.  So anyone who
wants to defend set theory as foundations has to respond it's just a code.
And the code isn't unique.  If real numbers are (coded as) equivalence
classes of Cauchy sequences, then elements of $\sqrt{2}$ are
rational-valued Cauchy sequences converging to $\sqrt{2}$.  If real numbers
are (coded as) Dedekind cuts, then the elements of $\sqrt{2}$ are the
sets $\set{ x \in \rats : x^2 < 2 }$ and $\set{ x \in \rats : x^2 > 2 }$.
But many other codings could be devised.
To do real analysis one does not want to know any of this.  Rather one
treats the real numbers as a complete (Cauchy complete or Dedekind complete,
whichever you want) Archimedian ordered field.  Note that that is what
the Cauchy and Dedekind reals $\real_c$ and $\real_d$ are
in HoTT.\footnote{More precisely, $\real_d$ is a mere set
that is an Archimedean ordered field \citepalias[Theorem~11.2.8]{hott-book}
that is both Cauchy and Dedekind complete \citepalias[Theorem~11.2.12 and
Corollary~11.2.16]{hott-book}, and $\real_c$ is a mere set
\citepalias[Theorem~11.3.9]{hott-book} that is an Archimedean ordered field
\citepalias[Theorem~11.3.48]{hott-book} that is Cauchy complete
\citepalias[Theorem~11.3.49]{hott-book}.

What ordered field means \citepalias[Definition~11.2.7]{hott-book} is different
intuitionistically and classically.
Classically, one has trichotomy: for all $x$ and $y$ one of $x = y$, $x < y$,
and $x > y$ holds.
Intuitionistically, one only has a weak total order: $x < y$ implies for all
$z$ either $z < y$ or $z > x$.}
If we have $x : \real$, then $x$ is just a point in that space.
It is not a type that has terms itself (unlike set theory).

So the reason why formalization in logic and set theory never caught on
with mathematicians is that it is profoundly unnatural.  Mathematicians
want types (whether or not they want formal type theory).  MLTT and HoTT
have the types and also have only operations that mathematicians find familiar
(functions and ordered pairs).  They were designed by mathematicians to be
useful for doing mathematics.  It remains to be seen how successful they
will be at that.

\section{Confession}

I should confess that, despite my enthusiasm about the methodology discussed
in this document, I have never used any of it in published research.

I have done nonstandard analysis
\citep{geyer-nsa,de-andrade-geyer,geyer-de-andrade},
and there were some surprising connections between that research and
both constructive mathematics \citep{bishop} and finitely-additive
probability theory,
which is another area of non-mainstream mathematics.\footnote{The
connections were that we were using ``radically elementary probability theory''
following \citet{nelson}
in which all probability spaces were finite although perhaps nonstandard
and this makes Kolmogorov's axiom of countable additivity vacuous.
Moreover, Birkhoff's ergodic theorem is not provable either constructively
\citep{bishop} nor in Nelson-style probability theory.}

I have also been intrigued by synthetic mathematics.  I am old enough to
have had a year of Euclidean geometry in high school.  Also long
ago in FIG we read \citet{bell}, a book about synthetic differential geometry.

I did have a PhD advisee that did a little bit of computer theorem proving.
It turned out to be too time-consuming for us and the student switched
from using proof assistants \citep[Chapter~6]{johnson-thesis} to proving
theorems the old-fashioned
way \citep[the rest of the thesis and][]{johnson-geyer}.
Perhaps we were using the wrong proof assistants
(not based on HoTT).\footnote{We used first ACL2 and then HOL Light.}
The main reason I am so intrigued by HoTT is it seems so much closer
to conventional mathematics than other formalizations.

But I have not actually used constructive mathematics or category theory
for research or teaching (except for one week of category theory in
a special topics course).

\section{Conclusion}

I do not wish to draw any super-strong conclusions from this.
I am satisfied to say that philosophers should be as intrigued about this
stuff as I am.  That having been said, I will not be wishy-washy in what
follows (hyper-qualification is boring writing).

Writing this has converted me to the philosophy that equality does not belong
in the logic.  I was long ago converted to the up-to-isomorphism view by all
of the mathematics I know.  In probability and statistics we treat random
variables as equal if they are equal in distribution (isomorphic in the
category where the objects are probability measures and the morphisms are
measurable functions).  I was amazed at how much simpler the explanation
of the first isomorphism theorem of group theory
(Section~\ref{sec:abstract-algebra} above) when I didn't have to use
equivalence classes and could just redefine what equality means for each
group.

All of real analysis, functional analysis, probability theory, and mathematical
statistics has taught me that the function concept is way more fundamental than
the set concept, so ETCS and category theory are on the right track here.
I say this, even though a probability measure is a function
$\mathcal{A} \to \real$ where $\mathcal{A}$ is a family of sets
(a sigma-algebra, footnote~\ref{foot:sigma-algebra}).
One can do probability theory without the set theory so functions are the
fundamental concept \citep{whittle}.\footnote{This has the weird effect
that probability theory becomes \emph{expectation theory}, probability
is just a special case of expectation and expectation (of random variables)
are what the axioms are about, when we rewrite them following \citet{whittle}.
Expectation is a function that maps random variables to real numbers.}
This accounts for why I am intrigued by ECTS and category theory.

As for MLTT and HoTT, I would like to get my feet wet with that (but perhaps
am too old with too many other projects to work on).  Constructive probability
theory, done like in \citet{bishop} with countable choice
so $\real_c \simeq \real_d$, or done like in HoTT so we must choose whether
we want $\real_c$ or $\real_d$ for expected values of random variables,
is doable (there is a lot in \citet{bishop} and more
in \citet{bishop-bridges}).

\begin{thebibliography}{}

\bibitem[Aluffi(2009)]{aluffi}
Aluffi, P. (2009).
\newblock \emph{Algebra, Chapter 0}.
\newblock American Mathematical Society, Providence, RI.

\bibitem[Aschbacher(2004)]{aschbacher}
Aschbacher, M. (2004).
\newblock The status of the classification of the finite simple groups.
\newblock \emph{Notices of the American Mathematical Society}, \textbf{51},
    736--740.

\bibitem[Bauer(2017)]{bauer-constructive}
Bauer, A. (2017).
\newblock Five stages of accepting constructive mathematics.
\newblock \emph{Bulletin (New Series) of the American Mathematical Society},
    \textbf{54}, 481--498.
\newblock \doi{10.1090/bull/1556}.

\bibitem[Bauer(2019)]{bauer-hott}
Bauer, A. (2019).
\newblock Introduction to homotopy type theory.
\newblock A doctoral course on homotopy theory and homotopy type theory
    given by Andrej Bauer and Jaka Smrekar at the Faculty of Mathematics
    and Physics, University of Ljubljana, in the Spring of 2019.
\newblock The lectures on homotopy type theory by Bauer were recorded
    and are on \href{https://www.youtube.com/playlist?list=PL-47DDuiZOMCRDiXDZ1fI0TFLQgQqOdDa}{YouTube}.
\newblock The course notes for those lectures are on
    \href{https://github.com/andrejbauer/homotopy-type-theory-course}{GitHub}.

\bibitem[Beattie and Butzmann(2002)]{beattie-butzmann}
Beattie, R., and Butzmann, H.-P. (2002).
\newblock \emph{Convergence Structures and Applications to Functional Analysis}.
\newblock Kluwer Academic Publishers, Dordrecht, The Netherlands.

\bibitem[Bell(2008)]{bell}
Bell, J.~L. (2008).
\newblock \emph{A Primer of Infinitesimal Analysis}, 2nd ed.
\newblock Cambridge University Press.

\bibitem[Bezem, et al.(2019)Bezem, Coquand, and Huber]{bezem-coquand-huber}
Bezem, M., Coquand, T., and Huber, S. (2019).
\newblock The univalence axiom in cubical sets.
\newblock \emph{Journal of Automated Reasoning}, \textbf{63}, 159--171.
\newblock \doi{10.1007/s10817-018-9472-6}.

\bibitem[Birkhoff and Mac Lane(1941)]{birkhoff-maclane}
Birkhoff, G., and Mac Lane, S. (1941).
\newblock \emph{A Survey of Modern Algebra}.
\newblock Macmillan, New York.

\bibitem[Bishop(1967)]{bishop}
Bishop, E. (1967).
\newblock \emph{Foundations of Constructive Analysis}.
\newblock McGraw-Hill, New York.

\bibitem[Bishop and Bridges(1985)]{bishop-bridges}
Bishop, E., and Bridges, D. (1985).
\newblock \emph{Constructive Analysis}.
\newblock Springer-Verlag, Berlin.

\bibitem[de Andrade and Geyer(2011)]{de-andrade-geyer}
de Andrade, B.~B. and Geyer, C.~J. (2011).
\newblock Nonstandard central limit theorems for Markov chains.
\newblock \emph{International Journal of Uncertainty, Fuzziness and
    Knowledge-Based Systems}, \textbf{19}, 251--274.
\newblock \doi{10.1142/S021848851100699X}.

\bibitem[Dijkstra(1968)]{go-to-harmful}
Dijkstra, E.~W. (1968).
\newblock Letters to the editor: go to statement considered harmful.
\newblock \emph{Communications of the ACM}, \textbf{11}, 147--148.
\newblock \doi{10.1145/362929.362947}.

\bibitem[Eilenberg and Mac Lane(1945)]{eilenberg-maclane}
Eilenberg, S., and Mac Lane, S. (1945).
\newblock General theory of natural equivalences.
\newblock \emph{Transactions of the American Mathematical Society},
    \textbf{58}, 231--294.
\newblock \doi{10.1090/S0002-9947-1945-0013131-6}.

\bibitem[Geyer(2007)]{geyer-nsa}
Geyer, C.~J. (2007).
\newblock Radically elementary probability and statistics.
\newblock Technical Report No.~657.  School of Statistics,
    University of Minnesota.
\newblock \url{https://hdl.handle.net/11299/199667}.

\bibitem[Geyer and de Andrade(2015)]{geyer-de-andrade}
Geyer, C.~J. and de Andrade, B.~B. (2015).
\newblock Near equivalence on metric spaces and a nonstandard central
    limit theorem.
\newblock \emph{Journal of Logic \& Analysis}, \textbf{7}, 1--20.
\newblock \doi{10.4115/jla.2015.7.3}.

\bibitem[Goldblatt(1984)]{goldblatt}
Goldblatt, R. (1984).
\newblock \emph{Topoi: The categorical analysis of logic.}
\newblock Elsevier, Amsterdam.  Republished by Dover, 2006.

\bibitem[Halmos(1974)]{halmos}
Halmos, P.~R. (1974).
\newblock \emph{Naive Set Theory}.
\newblock Springer-Verlag, New York.
\newblock Originally published by Litton Educational Publishing, Inc. (1960).

\bibitem[Hofmann and Streicher(1998)]{hofmann-streicher}
Hofmann, M., and Streicher, T. (1998).
\newblock The groupoid interpretation of type theory.
\newblock In G. Sambin and J.~M. Smith, (eds.),
    \emph{Twenty-five years of constructive type theory (Venice, 1995)}.
\newblock Oxford Logic Guides, \textbf{36}, 83--111.
\newblock Oxford University Press, New York.
\newblock \doi{10.1093/oso/9780198501275.003.0008}.

\bibitem[Jech(2003)]{jech}
Jech, T. (2003).
\newblock \emph{Set Theory}, third millenium edition, revised and expanded.
\newblock Springer-Verlag, Berlin.

\bibitem[Johnson(2011)]{johnson-thesis}
Johnson, L.~T. (2011)
\newblock Geometric ergodicity of a random-walk Metropolis algorithm
    via variable transformation and computer aided reasoning in statistics.
\newblock Ph.D. thesis, University of Minnesota.
\newblock \url{https://hdl.handle.net/11299/113140}.

\bibitem[Johnson and Geyer(2012)]{johnson-geyer}
Johnson, L.~T. and Geyer, C.~J. (2012).
\newblock Variable transformation to obtain geometric ergodicity
    in the random-walk Metropolis algorithm.
\newblock \emph{Annals of Statistics}, \textbf{40}, 3050--3076.
\newblock Correction: (2013) \emph{Annals of Statistics}, \textbf{41}, 2698.
\newblock \doi{10.1214/12-AOS1048} and \doi{10.1214/13-AOS1146}.

\bibitem[Johnstone(1983)]{johnstone}
Johnstone, P.~T. (1983).
\newblock The point of pointless topology.
\newblock \emph{Bulletin of the American Mathematical Society},
    \textbf{8}, 41--53.
\newblock \doi{10.1090/S0273-0979-1983-15080-2}.

\bibitem[Kleiner(2007)]{kleiner}
Kleiner, I. (2007).
\newblock \emph{A History of Abstract Algebra}.
\newblock Boston: Birkh\"{a}user.

\bibitem[Ladyman and Presnell(2018)]{ladyman-presnell-foundations}
Ladyman, J., and Presnell, S. (2018).
\newblock Does homotopy type theory provide a foundation for mathematics?
\newblock \emph{The British Journal for the Philosophy of Science},
    \textbf{69}, 377--420.
\newblock \doi{10.1093/bjps/axw006}.

\bibitem[Ladyman and Presnell(2015)]{ladyman-presnell}
Ladyman, J., and Presnell, S. (2015).
\newblock Identity in homotopy type theory, Part I: The justification of
    path induction.
\newblock{Philosophia Mathematica}, \textbf{23}, 386--406.
\newblock \doi{10.1093/philmat/nkv014}.

\bibitem[Lang, 1993]{lang}
Lang, S. (1993).
\newblock \emph{Real and Functional Analysis}, third edition.
\newblock Springer-Verlag, New York.

\bibitem[Lawvere(1963)]{lawvere-thesis}
Lawvere, F. W. (1963).
\newblock Functorial Semantics of Algebraic Theories.
\newblock Ph.D. thesis, Columbia University.
\newblock Republished (2004) with an author’s comment and a supplement in
    \emph{Reprints in Theory and Applications of Categories}, \textbf{5},
    1--121.
\newblock \url{http://www.tac.mta.ca/tac/reprints/articles/5/tr5.pdf}.

\bibitem[Lawvere(1964)]{lawvere-etcs}
Lawvere, F. W. (1964).
\newblock An elementary theory of the category of sets.
\newblock \emph{Proceedings of the National Academy of Science of the U.S.A.},
    \textbf{52}, 1506--1511.
\newblock Expanded version (2005) with commentary by Colin McLarty and
    the author in \emph{Reprints in Theory and Applications of Categories},
    \textbf{11}, pp. 1--35.

\bibitem[Lawvere(1966)]{lawvere-cat-of-cat}
Lawvere, F. W. (1966).
\newblock The category of categories as a foundation for mathematics.
\newblock In S. Eilenberg, D. K. Harrison, S. MacLane, H. Röhrl (eds.),
    \emph{Proceedings of the Conference on Categorical Algebra - La Jolla
    1965}, 1--20.
\newblock Springer.
\newblock \doi{10.1007/978-3-642-99902-4}.

\bibitem[Leinster(2014)]{leinster}
Leinster, T. (2014).
\newblock Rethinking set theory.
\newblock \emph{American Mathematical Monthly}, \textbf{121}, 403--415.
\newblock \doi{10.4169/amer.math.monthly.121.05.403}.

\bibitem[Mac Lane(1986)]{maclane}
Mac Lane, S. (1986).
\newblock \emph{Mathematics Form and Function}.
\newblock Springer-Verlag, New York.

\bibitem[Mac Lane(1998)]{maclane-working}
Mac Lane, S. (1998).
\newblock \emph{Categories for the Working Mathematician}, second edition.
\newblock Springer-Verlag, New York.

\bibitem[Martin-L\"{o}f(1975)]{martin-lof-mltt}
Martin-L\"{o}f, P. (1975).
\newblock An intuitionistic theory of types: predicative part.
\newblock In: H. E. Rose, J. C. Shepherdson (eds.),
    \emph{Logic Colloquium ‘73, Proceedings of the Logic Colloquium,
    Studies in Logic and the Foundations of Mathematics}, \textbf{80}, 73--118.
\newblock Elsevier.
\newblock \doi{10.1016/S0049-237X(08)71945-1}.

\bibitem[Mazur(2008)]{mazur}
Mazur, B. (2008).
\newblock When is one thing equal to some other thing?
\newblock In \emph{Proof and Other Dilemmas: Mathematics and Philosophy},
B. Gold and R.~A. Simons (eds.)
\newblock Mathematical Association of America, Washington, DC.
\newblock pp.~221--242.
\newblock \doi{10.5948/UPO9781614445050.015}.

\bibitem[Nelson(1987)]{nelson}
Nelson, E. (1987).
\newblock \emph{Radically Elementary Probability Theory}.
\newblock Princeton University Press.

\bibitem[Pierce(2002)]{pierce}
Pierce, B.~C. (2002).
\newblock \emph{Types and Programming Languages}.
\newblock MIT Press.

\bibitem[Rijke(2025)]{rijke}
Rijke, E. (2025).
\newblock \emph{Introduction to Homotopy Type Theory}.
\newblock Cambridge University Press.
\newblock \doi{10.1017/9781108933568}.
\newblock A preprint is available at
    \href{https://arxiv.org/abs/2212.11082}{Ar$\chi$iv}.

\bibitem[The Univalent Foundations Program(2013)]{hott-book}
The Univalent Foundations Program (2013).
\newblock \emph{Homotopy Type Theory: Univalent Foundations of Mathematics}.
\newblock Institute for Advanced Study, Princeton, NJ.
\newblock \url{https://homotopytypetheory.org/book}.

\bibitem[van der Waerden(1930)]{van-der-waerden}
van der Waerden, B. L. (1930).
\newblock \emph{Moderne Algebra, Teil I} 1930, \emph{Teil II} 1931.
\newblock Springer-Verlag, Berlin.

\bibitem[Voevodsky(2015)]{voevodsky}
Voevodsky, V. (2015).
\newblock An experimental library of formalized mathematics based on
    univalent foundations.
\newblock \emph{Mathematical Structures in Computer Science}, \textbf{25},
    1278--1294.
\newblock \doi{10.1017/S0960129514000577}.

\bibitem[Whittle(2005)]{whittle}
Whittle, P. (2005).
\newblock \emph{Probability via Expectation}, 4th ed.
\newblock New York: Springer.

\bibitem[Wilson(2009)]{wilson}
Wilson, R. A. (2009).
\newblock \emph{The Finite Simple Groups}.
\newblock Springer-Verlag, London.

\end{thebibliography}

\end{document}

